% Introducción a HTML y CSS - Desarrollo de Software
% Compilar con: latexmk   (ver .latexmkrc; usa XeLaTeX y el tema de resources/)
\documentclass[aspectratio=43,10pt,t]{beamer}

\usepackage[spanish,es-noshorthands]{babel}
\usetheme{UTN}
\usepackage{deck}
\definecolor{ptNavyB}{HTML}{31394D}
\definecolor{ptGray}{HTML}{666666}
\setbeamersize{text margin left=0.017\paperwidth}
\AtBeginDocument{\setlength{\leftmargini}{1.65em}}

\hypersetup{pdftitle={Introducción a HTML y CSS - Desarrollo de Software},pdfauthor={UTN FRSF}}

\begin{document}

% 1
\begin{utnfondo}
  \utnnote[anchor=base west,text=white,font=\utnserif\fontsize{47}{52}\selectfont]{0.045}{0.284}{Introducción}
  \utnnote[anchor=base west,text=white,font=\utnserif\fontsize{47}{52}\selectfont]{0.085}{0.487}{a HTML y CSS}
  \utnsubtitulo{0.044}{0.625}{Desarrollo de Software}
  \utnpic[height=0.262\paperheight,keepaspectratio=false]{html5-css3}{0.601}{0.679}{0.349}
\end{utnfondo}

% 2
\begin{frame}\relax
\utnpic[height=0.419\paperheight,keepaspectratio=false]{html-js-css}{0.237}{0.196}{0.536}
\utnpic[height=0.419\paperheight,keepaspectratio=false]{html-js-css}{0.237}{0.196}{0.536}
\utnrect{ptNavyB}{1}{0.020}{0.731}{0.946}{0.201}
\utnnote[anchor=north west]{0.033}{0.742}{\begin{minipage}[t]{.92\paperwidth}\setbeamerfont{itemize/enumerate body}{size=\fontsize{10.7}{12.6}\selectfont}\fontsize{10.7}{12.6}\selectfont\begin{itemize}
\item \textbf{HTML} describe la estructura un documento.
\item \textbf{CSS} describe cómo debe visualizarse.
\item \textbf{JavaScript} describe cómo debe responder a interacciones con el usuario.
\end{itemize}\end{minipage}}
\end{frame}

% 3
\begin{utnfondo}
  \utnnote[anchor=base west,text=white,font=\utnserif\fontsize{51}{56}\selectfont]{0.044}{0.425}{HTML}
  \utnnote[anchor=base west,text=white,font=\utnserif\fontsize{15.2}{18}\selectfont]{0.049}{0.509}{HyperText Markup language}
  \utnpic[height=0.262\paperheight,keepaspectratio=false]{html5}{0.763}{0.688}{0.171}
\end{utnfondo}

% 4
\begin{frame}{¿Qué es HTML?}\relax
\begingroup
\setbeamerfont{itemize/enumerate body}{size=\fontsize{13.2}{16.4}\selectfont}
\setbeamerfont{itemize/enumerate subbody}{size=\fontsize{12.1}{15.1}\selectfont}
\setbeamerfont{itemize/enumerate subsubbody}{size=\fontsize{11.2}{13.9}\selectfont}
\begin{itemize}
  \item\relax HTML es un lenguaje de etiquetado para describir páginas web.
  \item\relax Contiene un conjunto de \textbf{tags} para definir porciones de una página que son interpretadas por un renderizador \flecha{} \textbf{Navegador WEB}
  \item\relax Cada elemento de un tag describe contenido diferente.
  \item\relax El propósito fundamental de HTML es proveer para un documento que se puede transferir a través de la WEB, una descripción de la semántica de su contenido y establecer la estructura del documento (jerarquía de los componentes).
\end{itemize}
\endgroup
\end{frame}

% 5
\begin{frame}{Estructura de un documento}\relax
\begingroup
\setbeamerfont{itemize/enumerate body}{size=\fontsize{12.7}{17.3}\selectfont}
\setbeamerfont{itemize/enumerate subbody}{size=\fontsize{11.7}{15.9}\selectfont}
\setbeamerfont{itemize/enumerate subsubbody}{size=\fontsize{10.7}{14.6}\selectfont}
\begin{itemize}
  \item\relax HTML define una \textbf{estructura principal} para un documento en formato de árbol
\end{itemize}
\endgroup
\utnpic[height=0.482\paperheight,keepaspectratio=false]{estructura-documento}{0.045}{0.391}{0.439}
\utnpic[height=0.548\paperheight,keepaspectratio=false]{head-body}{0.633}{0.355}{0.321}
\utnpic[height=0.482\paperheight,keepaspectratio=false]{estructura-documento}{0.045}{0.391}{0.439}
\utnpic[height=0.548\paperheight,keepaspectratio=false]{head-body}{0.633}{0.355}{0.321}
\utnline[ptBlue,line width=0.6pt,->]{0.705}{0.463}{0.484}{0.523}
\utnline[ptBlue,line width=0.6pt,->]{0.678}{0.740}{0.448}{0.716}
\end{frame}

% 6
\begin{frame}\relax
\utnnote[anchor=center,text=utnTitle,font=\utnserif\fontsize{12.9}{15.5}\selectfont]{0.4965}{0.1105}{Ejemplo HTML}

\utncode[style=htmlvs,x=0.091,y=0.183,w=0.817,h=0.783,bg=FFFFFF,fs=5.90,lh=8.52,padx=27.0,pady=0.3,num=14.2]{ejemplo.html}
\utnpic[height=0.0345\paperheight,keepaspectratio=false]{html5-icono}{0.098}{0.148}{0.0250}
\utnnote[anchor=west,font=\fontsize{5.3}{6}\selectfont]{0.129}{0.166}{\color[HTML]{616161}mi\_primer\_pagina\_web.html\ \ >\ \ …}
\end{frame}

% 7
\begin{frame}{Visualización}\relax
\utnpic[height=0.525\paperheight,keepaspectratio=false]{visualizacion}{0.017}{0.259}{0.967}
\end{frame}

% 8
\begin{frame}{DOCTYPE}\relax
\begin{minipage}[t]{0.438\paperwidth}
\begingroup
\setbeamerfont{itemize/enumerate body}{size=\fontsize{12.2}{16.6}\selectfont}
\setbeamerfont{itemize/enumerate subbody}{size=\fontsize{10.2}{13.8}\selectfont}
\setbeamerfont{itemize/enumerate subsubbody}{size=\fontsize{9.3}{12.7}\selectfont}
\begin{itemize}
  \item\relax El tag <!DOCTYPE> es el primer elemento que aparece en un documento HTML5 y precisamente se utiliza para indicar esta característica al navegador \textbf{el tipo de documento que va a abrir}.
  \item\relax Para un documento HTML se utiliza:
  \begin{itemize}
    \item\relax \textbf{\textcolor[HTML]{0F243E}{<!DOCTYPE html>.}}
  \end{itemize}
\end{itemize}
\endgroup
\end{minipage}
\utncode[style=htmlvs,x=0.500,y=0.196,w=0.789,h=0.756,bg=FFFFFF,fs=5.70,lh=8.22,padx=26.0,pady=0.3,num=13.7]{ejemplo.html}
\utnrect{ptGreen}{1.50}{0.511}{0.187}{0.208}{0.048}
\end{frame}

% 9
\begin{frame}{Tag HTML}\relax
\utnnote[anchor=west]{0.018}{0.500}{\begin{minipage}[t]{0.426\paperwidth}
\begingroup
\setbeamerfont{itemize/enumerate body}{size=\fontsize{13.2}{18.0}\selectfont}
\setbeamerfont{itemize/enumerate subbody}{size=\fontsize{12.1}{16.5}\selectfont}
\setbeamerfont{itemize/enumerate subsubbody}{size=\fontsize{11.2}{15.2}\selectfont}
\begin{itemize}
  \item\relax Es la raíz de todo el documento e indica el inicio y fin del código HTML en el documento.\\
\end{itemize}
\endgroup
\end{minipage}}
\utncode[style=htmlvs,x=0.500,y=0.196,w=0.789,h=0.756,bg=FFFFFF,fs=5.70,lh=8.22,padx=26.0,pady=0.3,num=13.7]{ejemplo.html}
\utnrect{ptGreen}{1.50}{0.512}{0.220}{0.208}{0.048}
\end{frame}

% 10
\begin{frame}{Tag HEAD}\relax
\begin{minipage}[t]{0.438\paperwidth}
\begingroup
\setbeamerfont{itemize/enumerate body}{size=\fontsize{13.2}{18.0}\selectfont}
\setbeamerfont{itemize/enumerate subbody}{size=\fontsize{12.1}{16.5}\selectfont}
\setbeamerfont{itemize/enumerate subsubbody}{size=\fontsize{11.2}{15.2}\selectfont}
\begin{itemize}
  \item\relax Posee la información que describe al documento (metadatos) y la información de recursos necesarios para mostrar el documento correctamente, como por ejemplo la ubicación de hojas de estilos o archivos de script.
\end{itemize}
\endgroup
\end{minipage}
\utncode[style=htmlvs,x=0.500,y=0.196,w=0.789,h=0.756,bg=FFFFFF,fs=5.70,lh=8.22,padx=26.0,pady=0.3,num=13.7]{ejemplo.html}
\utnrect{ptGreen}{1.50}{0.506}{0.250}{0.435}{0.129}
\end{frame}

% 11
\begin{frame}{Tag META}\relax
\begin{minipage}[t]{0.438\paperwidth}
\begingroup
\setbeamerfont{itemize/enumerate body}{size=\fontsize{13.2}{18.0}\selectfont}
\setbeamerfont{itemize/enumerate subbody}{size=\fontsize{12.2}{16.6}\selectfont}
\setbeamerfont{itemize/enumerate subsubbody}{size=\fontsize{11.2}{15.2}\selectfont}
\begin{itemize}
  \item\relax El elemento "meta" permite definir diferentes tipos de metadata en el documento.
  \item\relax Se coloca siempre en el "head" de una página.
  \item\relax Los atributos posibles son:
  \begin{itemize}
    \item\relax name, content, charset, http-equiv
  \end{itemize}
\end{itemize}
\endgroup
\end{minipage}
\utncode[style=htmlvs,x=0.500,y=0.196,w=0.789,h=0.756,bg=FFFFFF,fs=5.70,lh=8.22,padx=26.0,pady=0.3,num=13.7]{ejemplo.html}
\utnrect{ptGreen}{1.50}{0.506}{0.250}{0.435}{0.129}
\utnrect{red}{1.50}{0.522}{0.314}{0.402}{0.045}
\end{frame}

% 12
\begin{frame}{Tag META}\relax
\begingroup
\setbeamerfont{itemize/enumerate body}{size=\fontsize{13.2}{18.0}\selectfont}
\setbeamerfont{itemize/enumerate subbody}{size=\fontsize{12.2}{16.6}\selectfont}
\setbeamerfont{itemize/enumerate subsubbody}{size=\fontsize{11.2}{15.2}\selectfont}
\begin{itemize}
  \item\relax Otros atributos posibles de META Tags:
  \begin{itemize}
    \item\relax \textit{application name}: el nombre de la aplicación web de la cual la página actual es parte.
    \item\relax \textit{author}: el autor de la página.
    \item\relax \textit{description}: descripción de la página
    \item\relax \textit{generator}: nombre del software que generó la página si fuese generada así.
    \item\relax \textit{keywords}: lista separada por comas de strings que describen el contenido de la página.
    \item\relax \textit{open graph tags}: para compartir en redes sociales
  \end{itemize}
\end{itemize}
\endgroup
\end{frame}

% 13
\begin{frame}{Tag LINK}\relax
\begin{minipage}[t]{0.438\paperwidth}
\begingroup
\setbeamerfont{itemize/enumerate body}{size=\fontsize{9.2}{11.5}\selectfont}
\setbeamerfont{itemize/enumerate subbody}{size=\fontsize{8.5}{10.5}\selectfont}
\setbeamerfont{itemize/enumerate subsubbody}{size=\fontsize{7.8}{9.7}\selectfont}
\begin{itemize}
  \item\relax {\fontsize{9.2}{11.5}\selectfont El tag <link> define la relación entre un documento y un recurso externo.\par}
  \item\relax {\fontsize{8.5}{10.6}\selectfont "rel" especifica la relación entre el documento linkeado y el documento. Aquí puede ser de utilidad el atributo \textbf{“prefetch”} que un recurso se pueda cargar en los tiempos ociosos del navegador, lo que hace que se cargue más rápido la página en algunas ocasiones.\par}
  \item\relax {\fontsize{8.5}{10.6}\selectfont "href" especifica la ubicación del recurso que puede ser una ruta absoluta o relativa.\par}
  \item\relax {\fontsize{8.5}{10.6}\selectfont "type" especifica el tipo que soportará el recurso (json, xml, etc)\par}
\end{itemize}
\endgroup
\end{minipage}
\utncode[style=htmlvs,x=0.500,y=0.196,w=0.789,h=0.756,bg=FFFFFF,fs=5.70,lh=8.22,padx=26.0,pady=0.3,num=13.7]{ejemplo.html}
\end{frame}

% 14
\begin{frame}{Tag LINK}\relax
{\fontsize{13.2}{18.0}\selectfont ¿Como vamos a utilizarlo nosotros? Generalmente para importar nuestro archivo css con estilos\par}
\utncode[style=htmlk,x=0.122,y=0.392,w=0.755,h=0.430,bg=FFFFFF,fs=8.61,lh=11.26,padx=26.4,pady=2.7,num=10.3]{link-css.html}
\end{frame}

% 15
\begin{frame}{Tag SCRIPT}\relax
\begin{minipage}[t]{0.438\paperwidth}
\begingroup
\setbeamerfont{itemize/enumerate body}{size=\fontsize{12.3}{15.2}\selectfont}
\setbeamerfont{itemize/enumerate subbody}{size=\fontsize{11.3}{14.0}\selectfont}
\setbeamerfont{itemize/enumerate subsubbody}{size=\fontsize{10.4}{12.9}\selectfont}
\begin{itemize}
  \item\relax El elemento <\textit{script>} permite incluir secuencias de comando, ya sea definidos en el documento o referenciados a un archivo externo y controlar su ejecución. El tipo de \textit{script} más usado es \textit{JavaScript}.
  \item\relax Normalmente se lo utiliza dentro del elemento \textit{<head>}.
\end{itemize}
\endgroup
\end{minipage}
\utnnote[anchor=west]{0.506}{0.370}{\begin{minipage}[t]{0.426\paperwidth}
\begingroup
\setbeamerfont{itemize/enumerate body}{size=\fontsize{9.4}{11.7}\selectfont}
\setbeamerfont{itemize/enumerate subbody}{size=\fontsize{9.4}{11.7}\selectfont}
\setbeamerfont{itemize/enumerate subsubbody}{size=\fontsize{8.7}{10.8}\selectfont}
\begin{itemize}
  \item\relax Atributos:
  \begin{itemize}
    \item\relax \textit{\textbf{type}}: especifica el tipo de script, para ~~JavaScript puede ser omitido.
    \item\relax \textit{\textbf{src}}: especifica la URL para un archivo script externo.
    \item\relax \textit{\textbf{defer}}: le dice al navegador que no ejecute el script hasta que la página sea cargada completamente.
    \item\relax \textit{\textbf{async}}:especifica que el script será ejecutado asincrónicamente cuando esté disponible.
  \end{itemize}
\end{itemize}
\endgroup
\end{minipage}}
\utncode[style=htmlk,x=0.160,y=0.545,w=0.644,h=0.450,bg=FFFFFF,fs=5.72,lh=7.51,padx=17.6,pady=1.8,num=6.9]{script.html}
\utnrect{ptGreen}{1.50}{0.132}{0.658}{0.677}{0.198}
\end{frame}

% 16
\begin{frame}{Tag Body}\relax
\begingroup
\setbeamerfont{itemize/enumerate body}{size=\fontsize{11.5}{14.2}\selectfont}
\setbeamerfont{itemize/enumerate subbody}{size=\fontsize{10.6}{13.1}\selectfont}
\setbeamerfont{itemize/enumerate subsubbody}{size=\fontsize{9.7}{12.1}\selectfont}
\begin{itemize}
  \item\relax El elemento <body> de HTML representa el contenido de un documento HTML. 
  \item\relax Solo puede haber un elemento <body> en un documento.
  \item\relax Dentro de él, se define la estructura completa del documento a través de diferentes tags
\end{itemize}
\endgroup
\utnpic[height=0.444\paperheight,keepaspectratio=false]{layout-body}{0.409}{0.500}{0.452}
\utnpic[height=0.257\paperheight,keepaspectratio=false]{head-body}{0.132}{0.542}{0.151}
\utnpic[height=0.444\paperheight,keepaspectratio=false]{layout-body}{0.409}{0.500}{0.452}
\utnpic[height=0.257\paperheight,keepaspectratio=false]{head-body}{0.132}{0.542}{0.151}
\utnline[ptBeige,line width=1.2pt]{0.264}{0.635}{0.408}{0.502}
\utnline[ptBeige,line width=1.2pt]{0.269}{0.779}{0.414}{0.944}
\end{frame}

% 17
\begin{frame}{Tag BODY}\relax
\utncode[style=htmlvs,x=0.105,y=0.191,w=0.789,h=0.756,bg=FFFFFF,fs=5.70,lh=8.22,padx=26.0,pady=0.3,num=13.7]{ejemplo.html}
\utnpic[height=0.280\paperheight,keepaspectratio=false]{visualizacion}{0.549}{0.191}{0.516}
\utnrect{ptGreen}{1.50}{0.114}{0.374}{0.706}{0.544}
\end{frame}

% 18
\begin{frame}{¿Que tags se pueden usar en el body?}\relax
\utnnote[anchor=north west]{0.063}{0.270}{\begin{minipage}[t]{0.835\paperwidth}
\begingroup
\setbeamerfont{itemize/enumerate body}{size=\fontsize{13.2}{18.0}\selectfont}
\setbeamerfont{itemize/enumerate subbody}{size=\fontsize{12.1}{16.5}\selectfont}
\setbeamerfont{itemize/enumerate subsubbody}{size=\fontsize{11.2}{15.2}\selectfont}
\begin{enumerate}
  \item\relax Elementos que permiten definir secciones
  \item\relax Elementos que permiten agrupar contenido
  \item\relax Elementos que permiten dar “formato” al texto
  \item\relax Elementos que permiten definir datos tabulares
  \item\relax Elementos que permiten definir formularios
\end{enumerate}
\endgroup
\end{minipage}}
\utnpic[height=0.355\paperheight,keepaspectratio=false]{meme-programmer}{0.720}{0.601}{0.215}
\end{frame}

% 19
\begin{frame}{Elementos semánticos para definir secciones}\relax
\utnnote[anchor=north west]{0.035}{0.176}{\begingroup\setlength{\tabcolsep}{2pt}\renewcommand{\arraystretch}{1.0}\arrayrulewidth=0.6pt
\begin{tabular}{|>{\raggedright\arraybackslash}p{0.201\paperwidth}|>{\raggedright\arraybackslash}p{0.712\paperwidth}|}
\hline
{\fontsize{9.1}{11.0}\selectfont\raggedright \href{https://developer.mozilla.org/es/docs/Web/HTML/Elemento/section}{\uline{<section>}}} & {\fontsize{9.1}{11.0}\selectfont\raggedright Define una sección en un documento.} \\
\hline
{\fontsize{9.1}{11.0}\selectfont\raggedright \href{https://developer.mozilla.org/es/docs/Web/HTML/Elemento/nav}{\uline{<nav>}}} & {\fontsize{9.1}{11.0}\selectfont\raggedright Define una sección que solamente contiene enlaces de navegación} \\
\hline
{\fontsize{9.1}{11.0}\selectfont\raggedright \href{https://developer.mozilla.org/es/docs/Web/HTML/Elemento/article}{\uline{<article>}}} & {\fontsize{9.1}{11.0}\selectfont\raggedright Define contenido autónomo que podría existir independientemente del resto del contenido.} \\
\hline
{\fontsize{9.1}{11.0}\selectfont\raggedright \href{https://developer.mozilla.org/es/docs/Web/HTML/Elemento/aside}{\uline{<aside>}}} & {\fontsize{9.1}{11.0}\selectfont\raggedright Define algunos contenidos vagamente relacionados con el resto del contenido de la página. Si es removido, el contenido restante seguirá teniendo sentido} \\
\hline
{\fontsize{9.1}{11.0}\selectfont\raggedright \href{https://developer.mozilla.org/en-US/docs/HTML/Element/Heading_Elements}{\uline{<h1>,<h2>,<h3>,}}\href{https://developer.mozilla.org/en-US/docs/HTML/Element/Heading_Elements}{\uline{<h4>,<h5>,<h6>}}} & {\fontsize{9.1}{11.0}\selectfont\raggedright Los elemento de cabecera ~implementan seis niveles de cabeceras de documentos;~<h1>~es la de mayor y~<h6>~es la de menor importancia. Un elemento de cabecera describe brevemente el tema de la sección que introduce.} \\
\hline
{\fontsize{9.1}{11.0}\selectfont\raggedright \href{https://developer.mozilla.org/es/docs/Web/HTML/Elemento/header}{\uline{<header>}}} & {\fontsize{9.1}{11.0}\selectfont\raggedright Define la cabecera de una página o sección. Usualmente contiene un logotipo, el título del sitio Web y una tabla de navegación de contenidos.} \\
\hline
{\fontsize{9.1}{11.0}\selectfont\raggedright \href{https://developer.mozilla.org/es/docs/Web/HTML/Elemento/footer}{\uline{<footer>}}} & {\fontsize{9.1}{11.0}\selectfont\raggedright Define el pie de una página o sección. Usualmente contiene un mensaje de derechos de autoría, algunos enlaces a información legal o direcciones para dar información de retroalimentación.} \\
\hline
{\fontsize{9.1}{11.0}\selectfont\raggedright \href{https://developer.mozilla.org/es/docs/Web/HTML/Elemento/address}{\uline{<address>}}} & {\fontsize{9.1}{11.0}\selectfont\raggedright Define una sección que contiene información de contacto.} \\
\hline
{\fontsize{9.1}{11.0}\selectfont\raggedright \href{https://developer.mozilla.org/es/docs/Web/HTML/Elemento/main}{\uline{<main>}}} & {\fontsize{9.1}{11.0}\selectfont\raggedright Define el contenido principal o importante en el documento. Solamente existe un elemento~<main>~en el documento.} \\
\hline
\end{tabular}\endgroup}
\end{frame}

% 20
\begin{frame}{Elementos que permiten agrupar contenido}\relax
\utnnote[anchor=north west]{0.047}{0.174}{\begingroup\setlength{\tabcolsep}{2pt}\renewcommand{\arraystretch}{1.0}\arrayrulewidth=0.6pt
\begin{tabular}{|>{\raggedright\arraybackslash}p{0.169\paperwidth}|>{\raggedright\arraybackslash}p{0.705\paperwidth}|}
\hline
{\fontsize{9.1}{11.0}\selectfont\raggedright \href{https://developer.mozilla.org/es/docs/Web/HTML/Elemento/p}{\uline{<p>}}} & {\fontsize{9.1}{11.0}\selectfont\raggedright Define una parte que debe mostrarse como un párrafo.} \\
\hline
{\fontsize{9.1}{11.0}\selectfont\raggedright \href{https://developer.mozilla.org/es/docs/Web/HTML/Elemento/hr}{\uline{<hr>}}} & {\fontsize{9.1}{11.0}\selectfont\raggedright Representa un quiebre temático entre parrafos de una sección o articulo o cualquier contenido.} \\
\hline
{\fontsize{9.1}{11.0}\selectfont\raggedright \href{https://developer.mozilla.org/es/docs/Web/HTML/Elemento/pre}{\uline{<pre>}}} & {\fontsize{9.1}{11.0}\selectfont\raggedright Indica que su contenido esta preformateado y que este formato debe ser preservado.} \\
\hline
{\fontsize{9.1}{11.0}\selectfont\raggedright \href{https://developer.mozilla.org/es/docs/Web/HTML/Elemento/blockquote}{\uline{<blockquote>}}} & {\fontsize{9.1}{11.0}\selectfont\raggedright Representa una contenido citado desde otra fuente.} \\
\hline
{\fontsize{9.1}{11.0}\selectfont\raggedright \href{https://developer.mozilla.org/es/docs/Web/HTML/Elemento/ol}{\uline{<ol>}}} & {\fontsize{9.1}{11.0}\selectfont\raggedright Define una lista ordenada de artículos.} \\
\hline
{\fontsize{9.1}{11.0}\selectfont\raggedright \href{https://developer.mozilla.org/es/docs/Web/HTML/Elemento/ul}{\uline{<ul>}}} & {\fontsize{9.1}{11.0}\selectfont\raggedright Define una lista de artículos sin orden. \href{http://www.mclibre.org/consultar/htmlcss/html/html-listas.html}{\uline{http:/\allowbreak{}/\allowbreak{}www.mclibre.org/\allowbreak{}consultar/\allowbreak{}htmlcss/\allowbreak{}html/\allowbreak{}html-listas.html}}} \\
\hline
{\fontsize{9.1}{11.0}\selectfont\raggedright \href{https://developer.mozilla.org/es/docs/Web/HTML/Elemento/li}{\uline{<li>}}} & {\fontsize{9.1}{11.0}\selectfont\raggedright Define un artículo de una lista ennumerada.} \\
\hline
{\fontsize{9.1}{11.0}\selectfont\raggedright \href{https://developer.mozilla.org/es/docs/Web/HTML/Elemento/dl}{\uline{<dl>}}} & {\fontsize{9.1}{11.0}\selectfont\raggedright Define una lista de definiciones, es decir, una lista de términos y sus definiciones asociadas.} \\
\hline
{\fontsize{9.1}{11.0}\selectfont\raggedright \href{https://developer.mozilla.org/es/docs/Web/HTML/Elemento/dt}{\uline{<dt>}}} & {\fontsize{9.1}{11.0}\selectfont\raggedright Representa un término definido por el siguiente~<dd>.} \\
\hline
{\fontsize{9.1}{11.0}\selectfont\raggedright \href{https://developer.mozilla.org/es/docs/Web/HTML/Elemento/dd}{\uline{<dd>}}} & {\fontsize{9.1}{11.0}\selectfont\raggedright Representa la definición de los terminos listados antes que él.} \\
\hline
{\fontsize{9.1}{11.0}\selectfont\raggedright \href{https://developer.mozilla.org/es/docs/Web/HTML/Elemento/figure}{\uline{<figure>}}} & {\fontsize{9.1}{11.0}\selectfont\raggedright Representa una figura ilustrada como parte ~del documento.} \\
\hline
{\fontsize{9.1}{11.0}\selectfont\raggedright \href{https://developer.mozilla.org/es/docs/Web/HTML/Elemento/figcaption}{\uline{<figcaption>}}} & {\fontsize{9.1}{11.0}\selectfont\raggedright Representa la leyenda de una figura.} \\
\hline
{\fontsize{9.1}{11.0}\selectfont\raggedright \href{https://developer.mozilla.org/es/docs/Web/HTML/Elemento/div}{\uline{<div>}}} & {\fontsize{9.1}{11.0}\selectfont\raggedright Representa un contenedor genérico sin ningún significado especial.} \\
\hline
\end{tabular}\endgroup}
\end{frame}

% 21
\begin{frame}{Semántica a nivel de Texto}\relax
\utnnote[anchor=north west,align=left]{0.292}{0.956}{\fontsize{9.1}{10.9}\selectfont \href{http://www.mclibre.org/consultar/htmlcss/html/html-texto.html}{\uline{http:/\allowbreak{}/\allowbreak{}www.mclibre.org/\allowbreak{}consultar/\allowbreak{}htmlcss/\allowbreak{}html/\allowbreak{}html-texto.html}}}
\utnnote[anchor=north west]{0.018}{0.175}{\begingroup\setlength{\tabcolsep}{2pt}\renewcommand{\arraystretch}{1.0}\arrayrulewidth=0.6pt
\begin{tabular}{|>{\raggedright\arraybackslash}p{0.122\paperwidth}|>{\raggedright\arraybackslash}p{0.807\paperwidth}|}
\hline
{\fontsize{8.1}{9.7}\selectfont\raggedright \href{https://developer.mozilla.org/es/docs/Web/HTML/Elemento/a}{\uline{<a>}}} & {\fontsize{8.1}{9.7}\selectfont\raggedright Representa un~hiperenlace~, enlazando a otro recurso. (\href{http://www.mclibre.org/consultar/htmlcss/html/html-enlaces.html}{\uline{http:/\allowbreak{}/\allowbreak{}www.mclibre.org/\allowbreak{}consultar/\allowbreak{}htmlcss/\allowbreak{}html/\allowbreak{}html-enlaces.html}} ) } \\
\hline
{\fontsize{8.1}{9.7}\selectfont\raggedright \href{https://developer.mozilla.org/es/docs/Web/HTML/Elemento/em}{\uline{<em>}}} & {\fontsize{8.1}{9.7}\selectfont\raggedright Representa un texto~enfatizado~, como un acento de intensidad.} \\
\hline
{\fontsize{8.1}{9.7}\selectfont\raggedright \href{https://developer.mozilla.org/es/docs/Web/HTML/Elemento/strong}{\uline{<strong>}}} & {\fontsize{8.1}{9.7}\selectfont\raggedright Representa un texto especialmente~importante~.} \\
\hline
{\fontsize{8.1}{9.7}\selectfont\raggedright \href{https://developer.mozilla.org/es/docs/Web/HTML/Elemento/small}{\uline{<small>}}} & {\fontsize{8.1}{9.7}\selectfont\raggedright Representa un~comentario aparte~, es decir, textos como un descargo de responsabilidad o una nota de derechos de autoría, que no son esenciales para la comprensión del documento.} \\
\hline
{\fontsize{8.1}{9.7}\selectfont\raggedright \href{https://developer.mozilla.org/es/docs/Web/HTML/Elemento/time}{\uline{<time>}}} & {\fontsize{8.1}{9.7}\selectfont\raggedright Representa un valor de~fecha~y~hora; el equivalente legible por máquina puede ser representado en el atributo~datetime.} \\
\hline
{\fontsize{8.1}{9.7}\selectfont\raggedright \href{https://developer.mozilla.org/es/docs/Web/HTML/Elemento/code}{\uline{<code>}}} & {\fontsize{8.1}{9.7}\selectfont\raggedright Representa un~código de~ordenador~.} \\
\hline
{\fontsize{8.1}{9.7}\selectfont\raggedright \href{https://developer.mozilla.org/es/docs/Web/HTML/Elemento/sub}{\uline{<sub>}},\href{https://developer.mozilla.org/es/docs/Web/HTML/Elemento/sup}{\uline{<sup>}}} & {\fontsize{8.1}{9.7}\selectfont\raggedright Representan un~subíndice~y un~superíndice,~respectivamente.} \\
\hline
{\fontsize{8.1}{9.7}\selectfont\raggedright \href{https://developer.mozilla.org/es/docs/Web/HTML/Elemento/i}{\uline{<i>}}} & {\fontsize{8.1}{9.7}\selectfont\raggedright Representa un texto en una voz o estado de ánimo~alterno, o por lo menos de diferente calidad, como una designación taxonómica, un término técnico, una frase idiomática, un pensamiento o el nombre de un barco.} \\
\hline
{\fontsize{8.1}{9.7}\selectfont\raggedright \href{https://developer.mozilla.org/es/docs/Web/HTML/Elemento/b}{\uline{<b>}}} & {\fontsize{8.1}{9.7}\selectfont\raggedright Representa un texto hacia el cual se llama la atención para~propósitos utilitaros.~No confiere ninguna importancia adicional y no implica una voz alterna.} \\
\hline
{\fontsize{8.1}{9.7}\selectfont\raggedright \href{https://developer.mozilla.org/es/docs/Web/HTML/Elemento/u}{\uline{<u>}}} & {\fontsize{8.1}{9.7}\selectfont\raggedright Representa una anotación no textual~sin-articular, como etiquetar un texto como mal escrito o etiquetar un nombre propio en texto en Chino.} \\
\hline
{\fontsize{8.1}{9.7}\selectfont\raggedright \href{https://developer.mozilla.org/es/docs/Web/HTML/Elemento/span}{\uline{<span>}}} & {\fontsize{8.1}{9.7}\selectfont\raggedright Representa texto sin un significado específico. ~Este debe ser usado cuando~ningún otro~elemento semántico le confiere un significado adecuado, en cuyo caso, provendrá de atributos globales como~class,~lang, o~dir.} \\
\hline
{\fontsize{8.1}{9.7}\selectfont\raggedright \href{https://developer.mozilla.org/es/docs/Web/HTML/Elemento/br}{\uline{<br>}}} & {\fontsize{8.1}{9.7}\selectfont\raggedright Representa un~salto de línea.} \\
\hline
\end{tabular}\endgroup}
\end{frame}

% 22
\begin{frame}{Datos tabulares}\relax
\utnnote[anchor=north west]{0.022}{0.173}{\begingroup\setlength{\tabcolsep}{2pt}\renewcommand{\arraystretch}{1.0}\arrayrulewidth=0.6pt
\begin{tabular}{|>{\raggedright\arraybackslash}p{0.122\paperwidth}|>{\raggedright\arraybackslash}p{0.815\paperwidth}|}
\hline
{\fontsize{10.2}{12.2}\selectfont\raggedright \href{https://developer.mozilla.org/es/docs/Web/HTML/Elemento/table}{\uline{<table>}}} & {\fontsize{12.2}{14.6}\selectfont\raggedright Representa~datos con más de una dimensión.} \\
\hline
{\fontsize{10.2}{12.2}\selectfont\raggedright \href{https://developer.mozilla.org/es/docs/Web/HTML/Elemento/caption}{\uline{<caption>}}} & {\fontsize{12.2}{14.6}\selectfont\raggedright Representa el~título de una tabla.} \\
\hline
{\fontsize{10.2}{12.2}\selectfont\raggedright \href{https://developer.mozilla.org/es/docs/Web/HTML/Elemento/colgroup}{\uline{<colgroup>}}} & {\fontsize{12.2}{14.6}\selectfont\raggedright Representa un~conjunto de una o más columnas~de una tabla.} \\
\hline
{\fontsize{10.2}{12.2}\selectfont\raggedright \href{https://developer.mozilla.org/es/docs/Web/HTML/Elemento/col}{\uline{<col>}}} & {\fontsize{12.2}{14.6}\selectfont\raggedright Representa una~columna~de una tabla.} \\
\hline
{\fontsize{10.2}{12.2}\selectfont\raggedright \href{https://developer.mozilla.org/es/docs/Web/HTML/Elemento/tbody}{\uline{<tbody>}}} & {\fontsize{12.2}{14.6}\selectfont\raggedright Representa el bloque de filas que describen los~datos contretos~de una tabla.} \\
\hline
{\fontsize{10.2}{12.2}\selectfont\raggedright \href{https://developer.mozilla.org/es/docs/Web/HTML/Elemento/thead}{\uline{<thead>}}} & {\fontsize{12.2}{14.6}\selectfont\raggedright Representa el bloque de filas que describen las~etiquetas de columna~de una tabla.} \\
\hline
{\fontsize{10.2}{12.2}\selectfont\raggedright \href{https://developer.mozilla.org/es/docs/Web/HTML/Elemento/tfoot}{\uline{<tfoot>}}} & {\fontsize{12.2}{14.6}\selectfont\raggedright Representa los bloques de filas que describen los~resúmenes de columna~de una tabla.} \\
\hline
{\fontsize{10.2}{12.2}\selectfont\raggedright \href{https://developer.mozilla.org/es/docs/Web/HTML/Elemento/tr}{\uline{<tr>}}} & {\fontsize{12.2}{14.6}\selectfont\raggedright Representa una~fila de celdas~en una tabla.} \\
\hline
{\fontsize{10.2}{12.2}\selectfont\raggedright \href{https://developer.mozilla.org/es/docs/Web/HTML/Elemento/td}{\uline{<td>}}} & {\fontsize{12.2}{14.6}\selectfont\raggedright Representa una~celda de datos~en una tabla.} \\
\hline
{\fontsize{10.2}{12.2}\selectfont\raggedright \href{https://developer.mozilla.org/es/docs/Web/HTML/Elemento/th}{\uline{<th>}}} & {\fontsize{12.2}{14.6}\selectfont\raggedright Representa una~celda encabezado~en una tabla.~} \\
\hline
\end{tabular}\endgroup}
\end{frame}

% 23
\begin{frame}{Elementos de formularios}\relax
\utnnote[anchor=north west,align=left]{0.063}{0.898}{\fontsize{9.1}{10.9}\selectfont \href{http://www.mclibre.org/consultar/htmlcss/html/html-formularios.html}{\uline{http:/\allowbreak{}/\allowbreak{}www.mclibre.org/\allowbreak{}consultar/\allowbreak{}ht\\mlcss/\allowbreak{}html/\allowbreak{}html-formularios.html}}}
\utnnote[anchor=north west]{0.020}{0.193}{\begingroup\setlength{\tabcolsep}{2pt}\renewcommand{\arraystretch}{1.0}\arrayrulewidth=0.6pt
\begin{tabular}{|>{\raggedright\arraybackslash}p{0.138\paperwidth}|>{\raggedright\arraybackslash}p{0.799\paperwidth}|}
\hline
{\fontsize{9.1}{11.0}\selectfont\raggedright \href{https://developer.mozilla.org/es/docs/Web/HTML/Elemento/form}{\uline{<form>}}} & {\fontsize{9.1}{11.0}\selectfont\raggedright Representa un~formulario, consistiendo de controles que puede ser enviado a un servidor para procesamiento.} \\
\hline
{\fontsize{9.1}{11.0}\selectfont\raggedright \href{https://developer.mozilla.org/es/docs/Web/HTML/Elemento/fieldset}{\uline{<fieldset>}}} & {\fontsize{9.1}{11.0}\selectfont\raggedright Representa un~conjunto de controles.} \\
\hline
{\fontsize{9.1}{11.0}\selectfont\raggedright \href{https://developer.mozilla.org/es/docs/Web/HTML/Elemento/legend}{\uline{<legend>}}} & {\fontsize{9.1}{11.0}\selectfont\raggedright Representa el~título~de un~<fieldset>.} \\
\hline
{\fontsize{9.1}{11.0}\selectfont\raggedright \href{https://developer.mozilla.org/es/docs/Web/HTML/Elemento/label}{\uline{<label>}}} & {\fontsize{9.1}{11.0}\selectfont\raggedright Representa el~título~de un control de formulario.} \\
\hline
{\fontsize{9.1}{11.0}\selectfont\raggedright \href{https://developer.mozilla.org/es/docs/Web/HTML/Elemento/input}{\uline{<input>}}} & {\fontsize{9.1}{11.0}\selectfont\raggedright Representa un~campo de datos escrito~que permite al usuario editar los datos.} \\
\hline
{\fontsize{9.1}{11.0}\selectfont\raggedright \href{https://developer.mozilla.org/es/docs/Web/HTML/Elemento/button}{\uline{<button>}}} & {\fontsize{9.1}{11.0}\selectfont\raggedright Representa un~botón~.} \\
\hline
{\fontsize{9.1}{11.0}\selectfont\raggedright \href{https://developer.mozilla.org/es/docs/Web/HTML/Elemento/select}{\uline{<select>}}} & {\fontsize{9.1}{11.0}\selectfont\raggedright Representa un control que permite la~selección entre un conjunto de opciones.} \\
\hline
{\fontsize{9.1}{11.0}\selectfont\raggedright \href{https://developer.mozilla.org/es/docs/Web/HTML/Elemento/option}{\uline{<option>}}} & {\fontsize{9.1}{11.0}\selectfont\raggedright Representa una~opción~en un elemento~<select>, o una sugerencia de un elemento~<datalist>.} \\
\hline
{\fontsize{9.1}{11.0}\selectfont\raggedright \href{https://developer.mozilla.org/es/docs/Web/HTML/Elemento/textarea}{\uline{<textarea>}}} & {\fontsize{9.1}{11.0}\selectfont\raggedright Representa un~control de edición de texto multi-línea.} \\
\hline
{\fontsize{9.1}{11.0}\selectfont\raggedright \href{https://developer.mozilla.org/es/docs/Web/HTML/Elemento/progress}{\uline{<progress>}}} & {\fontsize{9.1}{11.0}\selectfont\raggedright Representa el~progreso de finalización~de una tarea.} \\
\hline
{\fontsize{9.1}{11.0}\selectfont\raggedright \href{https://developer.mozilla.org/es/docs/Web/HTML/Elemento/meter}{\uline{<meter>}}} & {\fontsize{9.1}{11.0}\selectfont\raggedright Representa la~medida~escalar (o el valor fraccionario) dentro de un rango conocido.} \\
\hline
\end{tabular}\endgroup}
\end{frame}

% 24
\begin{frame}{Ejercicio}\relax
\utnpic[height=0.783\paperheight,keepaspectratio=false]{desotravel}{0.141}{0.192}{0.718}
\utnrect{ptGray}{0.38}{0.141}{0.192}{0.718}{0.783}
\end{frame}

% 25
\begin{frame}{Formularios}\relax
\begingroup
\setbeamerfont{itemize/enumerate body}{size=\fontsize{12.2}{15.1}\selectfont}
\setbeamerfont{itemize/enumerate subbody}{size=\fontsize{11.3}{14.0}\selectfont}
\setbeamerfont{itemize/enumerate subsubbody}{size=\fontsize{10.4}{12.9}\selectfont}
\begin{itemize}
  \item\relax Los formularios en HTML son elementos que usamos para capturar la entrada del usuario de diversas maneras.
  \item\relax Un formulario HTML es una sección de un documento que contiene contenido normal, código, elementos especiales llamados \href{http://html.conclase.net/w3c/html401-es/interact/forms.html\#form-controls}{\uline{\textit{controles}}} (checkboxes, radiobotones, menúes, etc.), y rótulos (labels) en esos controles. 
  \item\relax Los usuarios "completan" un formulario modificando sus controles 
  \begin{itemize}
    \item\relax introduciendo texto
    \item\relax seleccionando objetos de un menú
    \item\relax etc.
  \end{itemize}
\end{itemize}
\endgroup
\end{frame}

% 26
\begin{frame}{Botones}\relax
\begin{minipage}[t]{0.458\paperwidth}
\begingroup
\setbeamerfont{itemize/enumerate body}{size=\fontsize{11.1}{13.7}\selectfont}
\setbeamerfont{itemize/enumerate subbody}{size=\fontsize{10.2}{12.7}\selectfont}
\setbeamerfont{itemize/enumerate subsubbody}{size=\fontsize{9.4}{11.7}\selectfont}
\begin{itemize}
  \item\relax Botones
  \begin{itemize}
    \item\relax botones \textbf{de envío} (type=”submit”): envía un formulario. 
    \item\relax botones \textbf{de reinicialización} (type=”reset”): reinicia todos los controles a sus valores iniciales.
    \item\relax botones \textbf{pulsadores} (type=”button”): no tienen un comportamiento por defecto. Cada botón pulsador puede tener asociados scripts en el lado del cliente
  \end{itemize}
\end{itemize}
\endgroup
\end{minipage}
\utncode[style=htmlk,x=0.498,y=0.222,w=0.502,h=0.229,bg=FFFFFF,fs=3.91,lh=5.56,padx=12.9,pady=-0.1,num=4.7]{botones.html}
\utnpic[height=0.255\paperheight,keepaspectratio=false]{botones-resultado}{0.584}{0.555}{0.308}
\utnrect{ptGray}{0.75}{0.584}{0.555}{0.308}{0.255}
\end{frame}

% 27
\begin{frame}{Input}\relax
\begin{minipage}[t]{0.458\paperwidth}
\begingroup
\setbeamerfont{itemize/enumerate body}{size=\fontsize{8.1}{11.0}\selectfont}
\setbeamerfont{itemize/enumerate subbody}{size=\fontsize{8.1}{11.0}\selectfont}
\setbeamerfont{itemize/enumerate subsubbody}{size=\fontsize{7.5}{10.2}\selectfont}
\begin{itemize}
  \item\relax Input es muy útil para los formularios, ya que permite ingresar diferentes tipos de datos. Es posible utilizarlo como: 
  \begin{itemize}
    \item\relax entrada de texto
    \item\relax entrada de password
    \item\relax checkbox
    \item\relax radio
    \item\relax file
    \item\relax y otros…. 
  \end{itemize}
  \item\relax El tipo de control se define en el type, y su valor por defecto es "text".
  \item\relax El atributo name asigna el nombre de control.
  \item\relax El atributo value especifica el valor inicial del control. Es opcional excepto cuando el atributo type tenga el valor "radio" o "checkbox".
\end{itemize}
\endgroup
\end{minipage}
\utncode[style=htmlk,x=0.514,y=0.196,w=0.467,h=0.344,bg=FFFFFF,fs=3.65,lh=5.20,padx=12.9,pady=0.1,num=5.2]{input.html}
\utnpic[height=0.289\paperheight,keepaspectratio=false]{input-resultado}{0.545}{0.612}{0.403}
\end{frame}

% 28
\begin{frame}{Radio buttons}\relax
\begin{minipage}[t]{0.458\paperwidth}
\begingroup
\setbeamerfont{itemize/enumerate body}{size=\fontsize{10.2}{12.7}\selectfont}
\setbeamerfont{itemize/enumerate subbody}{size=\fontsize{9.4}{11.7}\selectfont}
\setbeamerfont{itemize/enumerate subsubbody}{size=\fontsize{8.7}{10.8}\selectfont}
\begin{itemize}
  \item\relax <input type=”radio” name=”opciones”...>
  \item\relax Los botones de radio son como los checkboxes, excepto en que cuando varios comparten el mismo nombre de control, son mutuamente exclusivos: 
  \begin{itemize}
    \item\relax cuando uno está "encendido", todos los demás con el mismo nombre se "apagan". 
  \end{itemize}
  \item\relax Si ninguno de los elementos \textcolor[HTML]{E50000}{<input>} de un conjunto de botones de radio especifica 'CHECKED', entonces el navegador debe marca el primer botón del conjunto inicialmente.
\end{itemize}
\endgroup
\end{minipage}
\utncode[style=htmlk,x=0.500,y=0.185,w=0.466,h=0.543,bg=FFFFFF,fs=3.55,lh=5.05,padx=12.5,pady=0.7,num=5.0]{radio.html}
\utnpic[height=0.278\paperheight,keepaspectratio=false]{radio-resultado}{0.586}{0.700}{0.294}
\utnrect{ptGreen}{1.50}{0.596}{0.390}{0.303}{0.205}
\end{frame}

% 29
\begin{frame}{Formularios: ¿Cómo los utilizamos?}\relax
\begin{minipage}[t]{0.424\paperwidth}
\begingroup
\setbeamerfont{itemize/enumerate body}{size=\fontsize{11.2}{15.2}\selectfont}
\setbeamerfont{itemize/enumerate subbody}{size=\fontsize{9.1}{12.4}\selectfont}
\setbeamerfont{itemize/enumerate subsubbody}{size=\fontsize{8.4}{11.4}\selectfont}
\begin{itemize}
  \item\relax El tag \textcolor[HTML]{E50000}{<form>} permite definir una función a ejecutar cuando se presione algún botón de envío (type submit)
  \begin{itemize}
    \item\relax ó se presiona enter dentro de algún input
    \item\relax ó un script llama al método requestSubmit() del form
  \end{itemize}
  \item\relax Esta función se define a través de la propiedad \textbf{onSubmit}
\end{itemize}
\endgroup
\end{minipage}
\utncode[style=htmlk,x=0.475,y=0.206,w=0.525,h=0.369,bg=FFFFFF,fs=4.08,lh=5.82,padx=11.9,pady=0.1,num=3.4]{onsubmit.html}
\utnpic[height=0.140\paperheight,keepaspectratio=false]{consola-submit}{0.410}{0.716}{0.575}
\end{frame}

% 30
\begin{frame}{Ejemplo}\relax
\utncode[style=htmlk,x=0.031,y=0.195,w=0.784,h=0.604,bg=FFFFFF,fs=6.99,lh=9.15,padx=34.2,pady=0.2,num=12.7,first=56]{form-action.html}
\utnpic[height=0.298\paperheight,keepaspectratio=false]{formulario-resultado}{0.614}{0.678}{0.362}
\end{frame}

% 31
\begin{frame}{Más sobre forms}\relax
\utnnote[anchor=center,font=\fontsize{8.1}{9.7}\selectfont]{0.500}{0.500}{\href{https://developer.mozilla.org/en-US/docs/Learn_web_development/Extensions/Forms}{\uline{https:/\allowbreak{}/\allowbreak{}developer.mozilla.org/\allowbreak{}en-US/\allowbreak{}docs/\allowbreak{}Learn\_web\_development/\allowbreak{}Extensions/\allowbreak{}Forms}}}
\end{frame}

% 32
\begin{utnfondo}
  \utnnote[anchor=base west,text=white,font=\utnserif\fontsize{51}{56}\selectfont]{0.044}{0.425}{CSS}
  \utnnote[anchor=base west,text=white,font=\utnserif\fontsize{15.2}{18}\selectfont]{0.049}{0.509}{Cascading Style Sheets}
  \utnpic[height=0.262\paperheight,keepaspectratio=false]{css3}{0.763}{0.688}{0.171}
\end{utnfondo}

% 33
\begin{frame}{CSS3}\relax
\begingroup
\setbeamerfont{itemize/enumerate body}{size=\fontsize{11.3}{14.0}\selectfont}
\setbeamerfont{itemize/enumerate subbody}{size=\fontsize{10.4}{12.9}\selectfont}
\setbeamerfont{itemize/enumerate subsubbody}{size=\fontsize{9.5}{11.8}\selectfont}
\begin{itemize}
  \item\relax CSS es el acrónimo de \textit{Cascading Style Sheets}.
  \item\relax Es un mecanismo que permite configurar mediante una sintaxis muy sencilla, las propiedades visuales de cualquier elemento HTML, como \textbf{fuentes, colores, bordes, márgenes, dimensión}, etc.
  \item\relax Permite de manera sencilla separar la estructura y el contenido de la forma en que se presenta.
  \item\relax Los estilos CSS se pueden definir en un archivo separado o en el mismo documento.
  \item\relax Una hoja de estilo se compone de una lista de reglas las cuales cada una posee uno o más~\textit{selectores}~y un~\textit{bloque de declaración}
  \item\relax CSS3 no es parte de la especificación HTML5, \textbf{pero es una tecnología clave cuando escribimos aplicaciones basadas en HTML5.}
\end{itemize}
\endgroup
\end{frame}

% 34
\begin{frame}{Definiendo una regla CSS}\relax
\begingroup
\setbeamerfont{itemize/enumerate body}{size=\fontsize{9.1}{10.7}\selectfont}
\setbeamerfont{itemize/enumerate subbody}{size=\fontsize{9.1}{10.7}\selectfont}
\setbeamerfont{itemize/enumerate subsubbody}{size=\fontsize{8.4}{9.8}\selectfont}
\begin{itemize}
  \item\relax Regla: cada uno de los estilos que componen una hoja de estilos CSS. 
  \begin{itemize}
    \item\relax Cada regla está compuesta de una parte de "selectores", un símbolo de "llave de apertura" (\{), otra parte denominada "declaración" y por último, un símbolo de "llave de cierre" (\}).
  \end{itemize}
  \item\relax Selector: indica el/los elemento/s a los que se aplica la regla CSS.
  \item\relax Declaración: especifica el/los estilo/s que se aplican a los elementos.
  \item\relax Propiedad: característica que se modifica en el elemento seleccionado.
  \item\relax Valor: establece el nuevo valor de la característica modificada en el elemento.
\end{itemize}
\endgroup
\utnpic[height=0.369\paperheight,keepaspectratio=false]{regla-css}{0.121}{0.603}{0.498}
\utnpic[height=0.369\paperheight,keepaspectratio=false]{html-css-casa}{0.694}{0.603}{0.220}
\end{frame}

% 35
\begin{frame}{Selectores}\relax
\begingroup
\setbeamerfont{itemize/enumerate body}{size=\fontsize{12.2}{15.8}\selectfont}
\setbeamerfont{itemize/enumerate subbody}{size=\fontsize{10.2}{13.2}\selectfont}
\setbeamerfont{itemize/enumerate subsubbody}{size=\fontsize{9.3}{12.1}\selectfont}
\begin{itemize}
  \item\relax Cuando queremos definir un estilo para un elemento HTML, lo que define CSS son reglas para seleccionar el elemento en cuestión que se denominan, \textbf{selectores.}
  \item\relax Existen distintos tipos de selectores.
  \item\relax \textbf{Selector universal: *}
  \begin{itemize}
    \item\relax {\fontsize{10.2}{13.2}\selectfont Selecciona todos los elementos de la página.\par}
    \item\relax {\fontsize{10.2}{13.2}\selectfont Por ejemplo eliminar el margen y el relleno de todos los elementos HTML:\par}
  \end{itemize}
  \begin{enumerate}
    \item\relax {\fontsize{9.1}{11.9}\selectfont \textit{\textbf{* \{}}\par}
    \item\relax {\fontsize{9.1}{11.9}\selectfont  ~\textit{\textbf{margin: 0;}}\par}
    \item\relax {\fontsize{9.1}{11.9}\selectfont  ~\textit{\textbf{padding: 0;}}\par}
    \item\relax {\fontsize{9.1}{11.9}\selectfont \textit{\textbf{\}}}\par}
  \end{enumerate}
  \begin{itemize}
    \item\relax {\fontsize{9.1}{11.9}\selectfont \textit{Su uso no es muy común de manera aislada, pero sí en combinación con otros selectores.}\par}
  \end{itemize}
\end{itemize}
\endgroup
\end{frame}

% 36
\begin{frame}{Ejemplo del Selector universal}\relax
\begingroup
\setbeamerfont{itemize/enumerate body}{size=\fontsize{10.2}{13.8}\selectfont}
\setbeamerfont{itemize/enumerate subbody}{size=\fontsize{9.3}{12.7}\selectfont}
\setbeamerfont{itemize/enumerate subsubbody}{size=\fontsize{8.6}{11.7}\selectfont}
\begin{itemize}
  \item\relax Sin utilizar CSS:
\end{itemize}
{\fontsize{9.5}{13.0}\selectfont\mbox{}\par}
{\fontsize{9.5}{13.0}\selectfont\mbox{}\par}
{\fontsize{9.5}{13.0}\selectfont\mbox{}\par}
{\fontsize{9.5}{13.0}\selectfont\mbox{}\par}
{\fontsize{9.5}{13.0}\selectfont\mbox{}\par}
{\fontsize{9.5}{13.0}\selectfont\mbox{}\par}
\begin{itemize}
  \item\relax Con el ejemplo anterior:
\end{itemize}
\endgroup
\utncode[style=css,x=0.045,y=0.659,w=0.200,h=0.246,bg=FFFFFF,fs=7.37,lh=10.71,padx=4.7,pady=0.8]{css-universal.css}
\utnrect{ptGray}{0.75}{0.045}{0.659}{0.200}{0.246}
\utnpic[height=0.347\paperheight,keepaspectratio=false]{universal-sin-css}{0.311}{0.196}{0.642}
\utnrect{ptGray}{0.75}{0.311}{0.196}{0.642}{0.347}
\utnpic[height=0.276\paperheight,keepaspectratio=false]{universal-con-css}{0.311}{0.644}{0.656}
\utnrect{ptGray}{0.75}{0.311}{0.644}{0.656}{0.276}
\end{frame}

% 37
\begin{frame}{Selectores}\relax
\begingroup
\setbeamerfont{itemize/enumerate body}{size=\fontsize{13.2}{18.0}\selectfont}
\setbeamerfont{itemize/enumerate subbody}{size=\fontsize{12.2}{16.6}\selectfont}
\setbeamerfont{itemize/enumerate subsubbody}{size=\fontsize{11.2}{15.2}\selectfont}
\begin{itemize}
  \item\relax \textbf{Selector de tipo o etiqueta:}
  \begin{itemize}
    \item\relax Selecciona todos los elementos de la página cuya etiqueta HTML coincide con el valor del selector. El siguiente ejemplo selecciona todos los párrafos de la página
    \item\relax Para utilizar este selector, solamente es necesario indicar el \textbf{nombre de una etiqueta HTML} (\textit{sin los caracteres < >}).
    \item\relax El siguiente ejemplo aplica diferentes estilos a los titulares y a los párrafos de una página HTML:
    \begin{itemize}
      \item\relax h1 \{ ~~color: red; \}
      \item\relax h2 \{ ~color: blue; \}
      \item\relax  p \{ ~~color: black; \}
    \end{itemize}
  \end{itemize}
\end{itemize}
\endgroup
\end{frame}

% 38
\begin{frame}{Ejemplo del Selector de tipo o etiqueta}\relax
\begingroup
\setbeamerfont{itemize/enumerate body}{size=\fontsize{10.2}{13.8}\selectfont}
\setbeamerfont{itemize/enumerate subbody}{size=\fontsize{9.3}{12.7}\selectfont}
\setbeamerfont{itemize/enumerate subsubbody}{size=\fontsize{8.6}{11.7}\selectfont}
\begin{itemize}
  \item\relax Sin utilizar CSS:
\end{itemize}
{\fontsize{9.4}{12.8}\selectfont\mbox{}\par}
{\fontsize{9.4}{12.8}\selectfont\mbox{}\par}
{\fontsize{9.4}{12.8}\selectfont\mbox{}\par}
{\fontsize{9.4}{12.8}\selectfont\mbox{}\par}
{\fontsize{9.4}{12.8}\selectfont\mbox{}\par}
{\fontsize{9.4}{12.8}\selectfont\mbox{}\par}
\begin{itemize}
  \item\relax Con el ejemplo anterior:
\end{itemize}
\endgroup
\utncode[style=css,x=0.032,y=0.630,w=0.180,h=0.331,bg=FFFFFF,fs=5.68,lh=8.11,padx=2.2,pady=-0.1]{css-tipo.css}
\utnrect{ptGray}{0.75}{0.032}{0.630}{0.180}{0.331}
\utnpic[height=0.347\paperheight,keepaspectratio=false]{tipo-sin-css}{0.301}{0.196}{0.642}
\utnrect{ptGray}{0.75}{0.301}{0.196}{0.642}{0.347}
\utnpic[height=0.354\paperheight,keepaspectratio=false]{tipo-con-css}{0.301}{0.619}{0.648}
\utnrect{ptGray}{0.75}{0.301}{0.619}{0.648}{0.354}
\end{frame}

% 39
\begin{frame}{Selectores}\relax
\begingroup
\setbeamerfont{itemize/enumerate body}{size=\fontsize{12.2}{15.1}\selectfont}
\setbeamerfont{itemize/enumerate subbody}{size=\fontsize{9.4}{11.6}\selectfont}
\setbeamerfont{itemize/enumerate subsubbody}{size=\fontsize{8.6}{10.7}\selectfont}
\begin{itemize}
  \item\relax {\fontsize{12.2}{15.1}\selectfont \textbf{Selectores de etiqueta agrupados}: si queremos aplicar el mismo estilo a más de una etiqueta, podemos indicarlo separándolos por \textbf{coma}:\par}
  \begin{enumerate}
    \item\relax \textbf{\textcolor[HTML]{0070C0}{h1, h2, h3}} \textcolor[HTML]{0070C0}{\{}
    \item\relax  ~\textcolor[HTML]{0070C0}{color: \#8A8E27;}
    \item\relax  ~\textcolor[HTML]{0070C0}{font-family: Arial, Helvetica, sans-serif;}
    \item\relax \textcolor[HTML]{0070C0}{\}}
  \end{enumerate}
  \item\relax {\fontsize{12.2}{15.1}\selectfont \textbf{Selector descendente:} \textbf{{\fontsize{10.2}{13.2}\selectfont s}}{\fontsize{10.2}{13.2}\selectfont elecciona los elementos que se encuentran dentro de otros elementos. Un elemento es descendiente de otro cuando se encuentra entre las etiquetas de apertura y de cierre del otro elemento.}\par}
  \item\relax {\fontsize{9.4}{11.6}\selectfont Ejemplo: aplicar un estilo a todos los elementos <span> que estén dentro de un elemento <p> (aunque no necesariamente de forma directa)\par}
\end{itemize}
\begin{enumerate}
  \item\relax {\fontsize{9.4}{11.6}\selectfont \textbf{\textcolor[HTML]{0070C0}{p span}} \textcolor[HTML]{0070C0}{\{ color: red; \}}\par}
\end{enumerate}
\endgroup
\end{frame}

% 40
\begin{frame}{Ejemplo del Selector de etiqueta agrupados}\relax
\begingroup
\setbeamerfont{itemize/enumerate body}{size=\fontsize{10.2}{13.8}\selectfont}
\setbeamerfont{itemize/enumerate subbody}{size=\fontsize{9.3}{12.7}\selectfont}
\setbeamerfont{itemize/enumerate subsubbody}{size=\fontsize{8.6}{11.7}\selectfont}
\begin{itemize}
  \item\relax Utilizando el HTML\\anterior:
\end{itemize}
{\fontsize{10.2}{13.8}\selectfont\mbox{}\par}
{\fontsize{10.2}{13.8}\selectfont\mbox{}\par}
{\fontsize{10.2}{13.8}\selectfont\mbox{}\par}
{\fontsize{10.2}{13.8}\selectfont\mbox{}\par}
{\fontsize{10.2}{13.8}\selectfont\mbox{}\par}
\begin{itemize}
  \item\relax Resultado:
\end{itemize}
\endgroup
\utncode[style=css,x=0.351,y=0.196,w=0.604,h=0.345,bg=FFFFFF,fs=7.99,lh=11.51,padx=3.7,pady=0.4]{css-agrupados.css}
\utnrect{ptGray}{0.75}{0.351}{0.196}{0.604}{0.345}
\utnpic[height=0.315\paperheight,keepaspectratio=false]{agrupados-resultado}{0.314}{0.617}{0.642}
\utnrect{ptGray}{0.75}{0.314}{0.617}{0.642}{0.315}
\utnnote[anchor=north west,align=center,text width=0.251\paperwidth,inner sep=2pt,fill=ptCallout,rounded corners=7pt,text=white,minimum height=0.240\paperheight,minimum width=0.263\paperwidth]{0.022}{0.650}{\fontsize{7.1}{8.5}\selectfont Importante la separación por COMAS ( , ) \\ Si no se confundirá con el otro selector}
\end{frame}

% 41
\begin{frame}{Ejemplo del Selector descendente}\relax
\begingroup
\setbeamerfont{itemize/enumerate body}{size=\fontsize{10.2}{13.8}\selectfont}
\setbeamerfont{itemize/enumerate subbody}{size=\fontsize{9.3}{12.7}\selectfont}
\setbeamerfont{itemize/enumerate subsubbody}{size=\fontsize{8.6}{11.7}\selectfont}
\begin{itemize}
  \item\relax Utilizando el CSS\\del ejemplo \\anterior:
\end{itemize}
{\fontsize{10.2}{13.8}\selectfont\mbox{}\par}
{\fontsize{10.2}{13.8}\selectfont\mbox{}\par}
{\fontsize{10.2}{13.8}\selectfont\mbox{}\par}
{\fontsize{10.2}{13.8}\selectfont\mbox{}\par}
{\fontsize{10.2}{13.8}\selectfont\mbox{}\par}
\begin{itemize}
  \item\relax Resultado:
\end{itemize}
\endgroup
\utncode[style=htmlsin,x=0.314,y=0.182,w=0.639,h=0.401,bg=FFFFFF,fs=6.37,lh=9.13,padx=4.2,pady=-0.2]{body-descendente.html}
\utnrect{ptGray}{0.75}{0.314}{0.182}{0.639}{0.401}
\utncode[style=css,x=0.032,y=0.687,w=0.289,h=0.246,bg=FFFFFF,fs=9.12,lh=12.97,padx=5.4,pady=0.5]{css-descendente.css}
\utnrect{ptGray}{0.75}{0.032}{0.687}{0.289}{0.246}
\utnpic[height=0.171\paperheight,keepaspectratio=false]{descendente-resultado}{0.352}{0.762}{0.621}
\utnrect{ptGray}{0.75}{0.352}{0.762}{0.621}{0.171}
\utnnote[anchor=north west,align=center,text width=0.496\paperwidth,inner sep=2pt,fill=ptCallout,rounded corners=7pt,text=white,minimum height=0.072\paperheight,minimum width=0.508\paperwidth]{0.389}{0.624}{\fontsize{7.1}{8.5}\selectfont Sin comas se interpreta como selector descendente}
\end{frame}

% 42
\begin{frame}{Selectores}\relax
\utnnote[anchor=north west]{0.037}{0.194}{\begin{minipage}[t]{0.907\paperwidth}
\begingroup
\setbeamerfont{itemize/enumerate body}{size=\fontsize{11.6}{14.4}\selectfont}
\setbeamerfont{itemize/enumerate subbody}{size=\fontsize{8.7}{10.8}\selectfont}
\setbeamerfont{itemize/enumerate subsubbody}{size=\fontsize{8.0}{9.9}\selectfont}
\begin{itemize}
  \item\relax \textbf{Selector de clase:}  podemos definir un estilo con un nombre de clase específico y luego aplicárselo al atributo “class” de cualquier elemento HTML.
  \item\relax Por ejemplo si definimos el estilo destacado como
  \begin{itemize}
    \item\relax \textbf{\textcolor[HTML]{244061}{.destacado \{ color: red; \}}}
  \end{itemize}
  \item\relax Luego podemos usarlo en un documento HTML
  \begin{enumerate}
    \item\relax <body>
    \item\relax  ~\textbf{\red{<p class="destacado">Lorem ipsum dolor sit amet...</p>}}
    \item\relax  ~<p>Nunc sed lacus et est adipiscing accumsan...</p>
    \item\relax </body>
  \end{enumerate}
\end{itemize}
\endgroup
\end{minipage}}
\utnnote[anchor=north west,align=center,text width=0.413\paperwidth,inner sep=2pt,fill=ptNavy,draw=ptLine,line width=1.0pt,rounded corners=8pt,text=white,minimum height=0.193\paperheight,minimum width=0.425\paperwidth]{0.031}{0.772}{\fontsize{8.1}{9.7}\selectfont El selector .destacado se interpreta como "cualquier elemento de la página cuyo atributo class sea igual a destacado"}
\utnnote[anchor=north west,align=center,text width=0.445\paperwidth,inner sep=2pt,fill=ptNavy,draw=ptLine,line width=1.0pt,rounded corners=8pt,text=white,minimum height=0.193\paperheight,minimum width=0.457\paperwidth]{0.500}{0.772}{\fontsize{8.1}{9.7}\selectfont Este selector permite en una misma página HTML que varios elementos diferentes utilicen el mismo valor en el atributo class}
\end{frame}

% 43
\begin{frame}{Ejemplo del Selector de clase}\relax
\begingroup
\setbeamerfont{itemize/enumerate body}{size=\fontsize{10.2}{13.8}\selectfont}
\setbeamerfont{itemize/enumerate subbody}{size=\fontsize{9.3}{12.7}\selectfont}
\setbeamerfont{itemize/enumerate subsubbody}{size=\fontsize{8.6}{11.7}\selectfont}
\begin{itemize}
  \item\relax HTML:
\end{itemize}
{\fontsize{10.2}{13.8}\selectfont\mbox{}\par}
{\fontsize{10.2}{13.8}\selectfont\mbox{}\par}
{\fontsize{10.2}{13.8}\selectfont\mbox{}\par}
{\fontsize{10.2}{13.8}\selectfont\mbox{}\par}
{\fontsize{10.2}{13.8}\selectfont\mbox{}\par}
{\fontsize{10.2}{13.8}\selectfont\mbox{}\par}
\begin{itemize}
  \item\relax Con el CSS del ejemplo anterior:
\end{itemize}
\endgroup
\utncode[style=htmlsin,x=0.171,y=0.272,w=0.809,h=0.220,bg=FFFFFF,fs=8.20,lh=11.77,padx=6.5,pady=-0.3]{body-clase.html}
\utnrect{ptGray}{0.75}{0.171}{0.272}{0.809}{0.220}
\utncode[style=css,x=0.039,y=0.652,w=0.336,h=0.283,bg=FFFFFF,fs=10.63,lh=15.28,padx=11.6,pady=-0.9]{css-clase.css}
\utnrect{ptGray}{0.75}{0.039}{0.652}{0.336}{0.283}
\utnpic[height=0.237\paperheight,keepaspectratio=false]{clase-resultado}{0.421}{0.664}{0.469}
\utnrect{ptGray}{0.75}{0.421}{0.664}{0.469}{0.237}
\end{frame}

% 44
\begin{frame}{Selectores}\relax
\utnnote[anchor=north west]{0.025}{0.180}{\begin{minipage}[t]{0.919\paperwidth}
\begingroup
\setbeamerfont{itemize/enumerate body}{size=\fontsize{13.2}{15.8}\selectfont}
\setbeamerfont{itemize/enumerate subbody}{size=\fontsize{12.2}{14.6}\selectfont}
\setbeamerfont{itemize/enumerate subsubbody}{size=\fontsize{11.2}{13.4}\selectfont}
\begin{itemize}
  \item\relax {\fontsize{13.2}{15.8}\selectfont Selector de clase combinado con etiqueta: \textit{{\fontsize{10.2}{13.2}\selectfont Aplicar una fuente “roja” solo a los elementos de tipo “p” que tienen un atributo class con valor “destacado”}}\par}
  \begin{itemize}
    \item\relax p.destacado \{ color: red \}
  \end{itemize}
  \item\relax {\fontsize{10.2}{12.2}\selectfont En cambio si queremos aplicarlo a cualquier elemento que tenga como valor del atributo class “destacado” y que esté dentro de un elemento <p>, en lugar de usar el punto dejamos un espacio.\par}
  \begin{itemize}
    \item[]
    \begin{itemize}
      \item\relax p destacado \{ color: red \}
    \end{itemize}
  \end{itemize}
  \item\relax {\fontsize{10.2}{12.2}\selectfont Combinar selectores\par}
\end{itemize}
\endgroup
\end{minipage}}
\utnline[ptBeige,line width=1.5pt]{0.381}{0.697}{0.446}{0.651}
\utnline[ptBeige,line width=1.5pt]{0.708}{0.746}{0.708}{0.816}
\utnnote[anchor=north west,align=left,text width=0.391\paperwidth,inner sep=2pt,fill=ptNavy,draw=ptLine,line width=1.0pt,text=white,minimum height=0.238\paperheight,minimum width=0.403\paperwidth]{0.027}{0.689}{\fontsize{8.1}{9.7}\selectfont 1.~.error \{ color: red; \} \\ 2.~.error.destacado \{ color: blue; \} \\ 3.~.destacado \{ font-size: 15px; \} \\ 4.~.especial ~\{ font-weight: bold; \}}
\utnnote[anchor=north west,align=center,text width=0.517\paperwidth,inner sep=2pt,fill=ptNavy,draw=ptLine,line width=1.0pt,rounded corners=8pt,text=white,minimum height=0.117\paperheight,minimum width=0.529\paperwidth]{0.452}{0.610}{\fontsize{9.1}{10.9}\selectfont <p class="especial destacado error">texto...</p>}
\utnnote[anchor=north west,align=center,text width=0.517\paperwidth,inner sep=2pt,fill=white,draw=ptNavy,line width=1.0pt,text=ptNavy,minimum height=0.133\paperheight,minimum width=0.529\paperwidth]{0.452}{0.835}{\fontsize{9.1}{10.9}\selectfont Se mostrará en azul dado que selecciona .error.destacado que \textbf{es más específico}}
\end{frame}

% 45
\begin{frame}{Ejemplo del Selectores combinados}\relax
\begingroup
\setbeamerfont{itemize/enumerate body}{size=\fontsize{10.2}{13.8}\selectfont}
\setbeamerfont{itemize/enumerate subbody}{size=\fontsize{9.3}{12.7}\selectfont}
\setbeamerfont{itemize/enumerate subsubbody}{size=\fontsize{8.6}{11.7}\selectfont}
\begin{itemize}
  \item\relax HTML:
\end{itemize}
{\fontsize{10.2}{13.8}\selectfont\mbox{}\par}
{\fontsize{10.2}{13.8}\selectfont\mbox{}\par}
{\fontsize{10.2}{13.8}\selectfont\mbox{}\par}
{\fontsize{10.2}{13.8}\selectfont\mbox{}\par}
\begin{itemize}
  \item\relax CSS:\\Resultado:
\end{itemize}
\endgroup
\utncode[style=htmlsin,x=0.047,y=0.266,w=0.921,h=0.160,bg=FFFFFF,fs=9.96,lh=14.18,padx=4.0,pady=0.3]{body-combinados.html}
\utnrect{ptGray}{0.75}{0.047}{0.266}{0.921}{0.160}
\utncode[style=css,x=0.175,y=0.471,w=0.262,h=0.492,bg=FFFFFF,fs=6.63,lh=9.60,padx=4.4,pady=-0.3]{css-combinados.css}
\utnrect{ptGray}{0.75}{0.175}{0.471}{0.262}{0.492}
\utnpic[height=0.141\paperheight,keepaspectratio=false]{combinados-resultado}{0.656}{0.639}{0.262}
\utnrect{ptGray}{0.75}{0.656}{0.639}{0.262}{0.141}
\end{frame}

% 46
\begin{frame}{Selectores}\relax
\begingroup
\setbeamerfont{itemize/enumerate body}{size=\fontsize{13.2}{18.0}\selectfont}
\setbeamerfont{itemize/enumerate subbody}{size=\fontsize{12.2}{16.6}\selectfont}
\setbeamerfont{itemize/enumerate subsubbody}{size=\fontsize{11.2}{15.2}\selectfont}
\begin{itemize}
  \item\relax \textbf{Selectores de ID}: En ocasiones, es necesario aplicar estilos CSS a \textbf{un único elemento de la página}. Aunque puede utilizarse un selector de clase para aplicar estilos a un único elemento, existe otro selector más eficiente en este caso.
  \item\relax El selector de ID permite seleccionar un elemento de la página a través del valor de su atributo id. 
  \item\relax La sintaxis es similar a la del selector de estilo, pero en lugar del prefijo “.” utiliza el prefijo “\#”.
  \begin{itemize}
    \item\relax \#destacado \{ color: red; \}
  \end{itemize}
\end{itemize}
\endgroup
\end{frame}

% 47
\begin{frame}{Ejemplo del Selectores combinados}\relax
\begingroup
\setbeamerfont{itemize/enumerate body}{size=\fontsize{10.2}{13.8}\selectfont}
\setbeamerfont{itemize/enumerate subbody}{size=\fontsize{9.3}{12.7}\selectfont}
\setbeamerfont{itemize/enumerate subsubbody}{size=\fontsize{8.6}{11.7}\selectfont}
\begin{itemize}
  \item\relax HTML:
\end{itemize}
{\fontsize{10.2}{13.8}\selectfont\mbox{}\par}
{\fontsize{10.2}{13.8}\selectfont\mbox{}\par}
{\fontsize{10.2}{13.8}\selectfont\mbox{}\par}
{\fontsize{10.2}{13.8}\selectfont\mbox{}\par}
\begin{itemize}
  \item\relax CSS:\\Resultado:
\end{itemize}
\endgroup
\utncode[style=htmlsin,x=0.185,y=0.218,w=0.749,h=0.170,bg=FFFFFF,fs=8.10,lh=11.57,padx=4.3,pady=-0.5]{body-id.html}
\utnrect{ptGray}{0.75}{0.185}{0.218}{0.749}{0.170}
\utncode[style=css,x=0.069,y=0.513,w=0.377,h=0.429,bg=FFFFFF,fs=9.90,lh=14.31,padx=4.7,pady=-0.5]{css-id.css}
\utnrect{ptGray}{0.75}{0.069}{0.513}{0.377}{0.429}
\utnpic[height=0.214\paperheight,keepaspectratio=false]{id-resultado}{0.557}{0.534}{0.377}
\utnrect{ptGray}{0.75}{0.557}{0.534}{0.377}{0.214}
\utnnote[anchor=north west,align=center,text width=0.391\paperwidth,inner sep=2pt,fill=ptCallout,rounded corners=7pt,text=white,minimum height=0.126\paperheight,minimum width=0.403\paperwidth]{0.547}{0.807}{\fontsize{7.1}{8.5}\selectfont Sólo podemos tener un único ID asociado, para estos casos se debe utilizar el selector de ~\textbf{class}}
\end{frame}

% 48
\begin{frame}{Selectores}\relax
\begingroup
\setbeamerfont{itemize/enumerate body}{size=\fontsize{13.2}{17.2}\selectfont}
\setbeamerfont{itemize/enumerate subbody}{size=\fontsize{12.2}{15.8}\selectfont}
\setbeamerfont{itemize/enumerate subsubbody}{size=\fontsize{11.2}{14.6}\selectfont}
\begin{itemize}
  \item\relax \textbf{Selector de hijos}: es ~similar al selector descendente, pero selecciona un elemento que es hijo \textbf{\textcolor[HTML]{17365D}{directo}} de otro elemento y se indica mediante el "signo de mayor que" (>):
  \item\relax Dado el estilo:
  \begin{itemize}
    \item\relax p > span \{ color: blue; \}
  \end{itemize}
  \item\relax Y dado el siguiente HTML
  \begin{itemize}
    \item\relax <p>\textcolor[HTML]{0070C0}{<span>Texto1</span>}</p>
    \item\relax <p><a href="\#">\red{<span>Texto2</span>}</a></p>
  \end{itemize}
  \item\relax Solo se aplicará a Texto1, que es el único elemento <span> hijo directo de un elemento <p> 
\end{itemize}
\endgroup
\utnpic[height=0.289\paperheight,keepaspectratio=false]{meme-css}{0.701}{0.415}{0.217}
\end{frame}

% 49
\begin{frame}{Ejemplo del Selector de hijo}\relax
\begingroup
\setbeamerfont{itemize/enumerate body}{size=\fontsize{10.2}{13.8}\selectfont}
\setbeamerfont{itemize/enumerate subbody}{size=\fontsize{9.3}{12.7}\selectfont}
\setbeamerfont{itemize/enumerate subsubbody}{size=\fontsize{8.6}{11.7}\selectfont}
\begin{itemize}
  \item\relax Utilizando el \\mismo HTML que\\para el selector\\descendente:
\end{itemize}
{\fontsize{10.2}{13.8}\selectfont\mbox{}\par}
{\fontsize{10.2}{13.8}\selectfont\mbox{}\par}
{\fontsize{10.2}{13.8}\selectfont\mbox{}\par}
\begin{itemize}
  \item\relax Resultado:
\end{itemize}
\endgroup
\utncode[style=htmlsin,x=0.314,y=0.182,w=0.639,h=0.401,bg=FFFFFF,fs=6.37,lh=9.13,padx=4.2,pady=-0.2]{body-descendente.html}
\utnrect{ptGray}{0.75}{0.314}{0.182}{0.639}{0.401}
\utncode[style=css,x=0.033,y=0.648,w=0.307,h=0.278,bg=FFFFFF,fs=10.22,lh=14.61,padx=4.7,pady=1.1]{css-hijo.css}
\utnrect{ptGray}{0.75}{0.033}{0.648}{0.307}{0.278}
\utnpic[height=0.181\paperheight,keepaspectratio=false]{hijo-resultado}{0.357}{0.736}{0.595}
\utnrect{ptGray}{0.75}{0.357}{0.736}{0.595}{0.181}
\utnrect{ptGold}{1.10}{0.105}{0.716}{0.025}{0.035}
\utnnote[anchor=north west,align=center,text width=0.496\paperwidth,inner sep=2pt,fill=ptCallout,rounded corners=7pt,text=white,minimum height=0.072\paperheight,minimum width=0.508\paperwidth]{0.389}{0.624}{\fontsize{7.1}{8.5}\selectfont Sólo afecta a los hijos directos (span dentro de p)}
\end{frame}

% 50
\begin{frame}{Unidades de Medida}\relax
\begingroup
\setbeamerfont{itemize/enumerate body}{size=\fontsize{13.2}{16.4}\selectfont}
\setbeamerfont{itemize/enumerate subbody}{size=\fontsize{10.2}{12.6}\selectfont}
\setbeamerfont{itemize/enumerate subsubbody}{size=\fontsize{9.3}{11.6}\selectfont}
\begin{itemize}
  \item\relax Mediante CSS se pueden definir la altura, el ancho, los márgenes y otras dimensiones de los elementos HTML.
  \item\relax Las unidades que se pueden utilizar son diversas y se pueden diferenciar las relativas de las \red{absolutas}
  \begin{itemize}
    \item\relax {\fontsize{10.2}{12.6}\selectfont \red{in, pulgadas ("inches", en inglés). Una pulgada equivale a 2.54 centímetros.}\par}
    \item\relax {\fontsize{10.2}{12.6}\selectfont \red{cm, centímetros.}\par}
    \item\relax {\fontsize{10.2}{12.6}\selectfont \red{mm, milímetros.}\par}
    \item\relax {\fontsize{9.1}{11.3}\selectfont \textbf{em}, \flechag{} relativa respecto del tamaño de letra del elemento.\par}
    \item\relax {\fontsize{9.1}{11.3}\selectfont \textbf{px, (píxel)} \flechag{} ~relativa respecto de la resolución de la pantalla del dispositivo en el que se visualiza la página HTML.\par}
    \item\relax {\fontsize{9.1}{11.3}\selectfont \textbf{vw, (view width)} \flechag{} equivale al \% del ancho de la ventana. Ej, un elemento con width: 50vw; ocupará la mitad del ancho de la pantalla visible.\par}
    \item\relax {\fontsize{9.1}{11.3}\selectfont \textbf{vh, (view height)} \flechag{} equivale al \% de la altura de la ventana. Ej, ejemplo, un elemento con height: 100vh; llenará completamente la altura de la pantalla.\par}
  \end{itemize}
\end{itemize}
\endgroup
\end{frame}

% 51
\begin{frame}{Colores}\relax
\begingroup
\setbeamerfont{itemize/enumerate body}{size=\fontsize{13.2}{18.0}\selectfont}
\setbeamerfont{itemize/enumerate subbody}{size=\fontsize{12.2}{16.6}\selectfont}
\setbeamerfont{itemize/enumerate subsubbody}{size=\fontsize{11.2}{15.2}\selectfont}
\begin{itemize}
  \item\relax Los colores en CSS se pueden indicar de cinco formas diferentes: palabras clave, colores del sistema, RGB hexadecimal, RGB numérico y RGB porcentual. Aunque el método más habitual es el del RGB hexadecimal.
  \item\relax Ejemplos
  \begin{itemize}
    \item\relax p \{ color: rgb(71, 98, 176); \}
    \item\relax p \{ color: \#4762B0; \}
  \end{itemize}
\end{itemize}
\endgroup
\utnpic[height=0.446\paperheight,keepaspectratio=false]{colores}{0.535}{0.490}{0.417}
\end{frame}

% 52
\begin{frame}{Buscar “rgb to hex” o similar}\relax
\utnpic[height=0.689\paperheight,keepaspectratio=false]{rgb-to-hex}{0.070}{0.209}{0.713}
\utnnote[anchor=north west,align=center,text width=0.154\paperwidth,inner sep=2pt,fill=ptCallout,rounded corners=14pt,text=white,minimum height=0.299\paperheight,minimum width=0.166\paperwidth]{0.803}{0.424}{\fontsize{7.1}{8.5}\selectfont Permite pasar de RGB a HEX\\o de\\HEX a RGB}
\end{frame}

% 53
\begin{frame}\relax
\utnnote[anchor=center,text=utnTitle,font=\utnserif\fontsize{14.0}{16.8}\selectfont]{0.4965}{0.1205}{Modelo de cajas}

\utnnote[anchor=north west]{0.010}{0.151}{\begin{minipage}[t]{0.942\paperwidth}
\begingroup
\setbeamerfont{itemize/enumerate body}{size=\fontsize{10.2}{13.2}\selectfont}
\setbeamerfont{itemize/enumerate subbody}{size=\fontsize{9.3}{12.1}\selectfont}
\setbeamerfont{itemize/enumerate subsubbody}{size=\fontsize{8.6}{11.2}\selectfont}
\begin{itemize}
  \item\relax CSS utiliza para representar todos los elementos un modelo de cajas que se compone de 6 partes. En principio no tienen visibilidad pero le permiten al desarrollador modificar para lograr la apariencia deseada.
\end{itemize}
\endgroup
\end{minipage}}
\utnpic[height=0.606\paperheight,keepaspectratio=false]{modelo-cajas}{0.500}{0.336}{0.483}
\utnnote[anchor=north west,align=left,text width=0.446\paperwidth,inner sep=2pt,fill=white,draw=ptNavy,line width=1.0pt,text=black,minimum height=0.657\paperheight,minimum width=0.458\paperwidth]{0.019}{0.311}{\fontsize{7.9}{9.5}\selectfont \textbf{content}: contenido HTML del elemento (palabras, una imagen, etc.) \\ \textbf{padding}: espacio libre opcional existente entre el contenido y el borde. \\ \textbf{border}: línea que encierra completamente el contenido y su relleno. \\ \textbf{background-image}: imagen que se muestra por detrás del contenido y el espacio de relleno. \\ \textbf{background-color}: color que se muestra por detrás del contenido y el espacio de relleno. \\ \textbf{margin}: ~separación opcional existente entre la caja y el resto de cajas adyacentes.}
\end{frame}

% 54
\begin{frame}{Modelo de cajas - Dimensiones}\relax
\utnnote[anchor=north west]{0.025}{0.178}{\begin{minipage}[t]{0.919\paperwidth}
\begingroup
\setbeamerfont{itemize/enumerate body}{size=\fontsize{13.2}{19.8}\selectfont}
\setbeamerfont{itemize/enumerate subbody}{size=\fontsize{12.2}{18.3}\selectfont}
\setbeamerfont{itemize/enumerate subsubbody}{size=\fontsize{11.2}{16.8}\selectfont}
\begin{itemize}
  \item\relax Para controlar los márgenes y el relleno de costado se definen los atributos
  \begin{itemize}
    \item\relax margin y padding
  \end{itemize}
  \item\relax Ancho y alto: {\fontsize{10.2}{13.2}\selectfont width} y {\fontsize{10.2}{13.2}\selectfont heigth}
  \item\relax Bordes:
  \begin{itemize}
    \item\relax border-width / border-color / border-style
  \end{itemize}
  \item\relax Ejemplo
\end{itemize}
\endgroup
\end{minipage}}
\utnpic[height=0.263\paperheight,keepaspectratio=false]{modelo-dimensiones}{0.479}{0.237}{0.477}
\utnnote[anchor=north west,align=left,text width=0.319\paperwidth,inner sep=2pt,fill=white,draw=teal,line width=1.0pt,text=black,minimum height=0.323\paperheight,minimum width=0.331\paperwidth]{0.028}{0.608}{\fontsize{8.6}{10.3}\selectfont \textbullet~~div \{ \\ \textbullet~~ ~width: 300px; \\ \textbullet~~ ~padding-left: ~50px; \\ \textbullet~~ ~padding-right: 50px; \\ \textbullet~~ ~margin-left: ~~30px; \\ \textbullet~~ ~margin-right: ~30px; \\ \textbullet~~ ~border: 10px solid black; \\ \textbullet~~\}}
\utnnote[anchor=north west,align=center,text width=0.594\paperwidth,inner sep=2pt,fill=ptNavy,draw=ptNavy,line width=1.0pt,text=white,minimum height=0.221\paperheight,minimum width=0.606\paperwidth]{0.374}{0.660}{\fontsize{8.4}{10.1}\selectfont El ancho del elemento en pantalla será la suma de la anchura original, los márgenes, los bordes y los rellenos: \\ 30px + 10px + 50px + 300px + 50px + 10px + 30px = 480 píxel}
\end{frame}

% 55
\begin{frame}{Ejemplo de <div>}\relax
\begingroup
\setbeamerfont{itemize/enumerate body}{size=\fontsize{10.2}{13.8}\selectfont}
\setbeamerfont{itemize/enumerate subbody}{size=\fontsize{9.3}{12.7}\selectfont}
\setbeamerfont{itemize/enumerate subsubbody}{size=\fontsize{8.6}{11.7}\selectfont}
\begin{itemize}
  \item\relax HTML:
\end{itemize}
{\fontsize{8.6}{11.7}\selectfont\mbox{}\par}
{\fontsize{8.6}{11.7}\selectfont\mbox{}\par}
{\fontsize{8.6}{11.7}\selectfont\mbox{}\par}
\begin{itemize}
  \item\relax CSS:
\end{itemize}
{\fontsize{8.6}{11.7}\selectfont\mbox{}\par}
{\fontsize{8.6}{11.7}\selectfont\mbox{}\par}
{\fontsize{8.6}{11.7}\selectfont\mbox{}\par}
{\fontsize{8.6}{11.7}\selectfont\mbox{}\par}
{\fontsize{8.6}{11.7}\selectfont\mbox{}\par}
\begin{itemize}
  \item\relax Resultado:
\end{itemize}
\endgroup
\utncode[style=htmlsin,x=0.216,y=0.196,w=0.632,h=0.128,bg=FFFFFF,fs=8.04,lh=11.55,padx=4.5,pady=1.8]{body-div.html}
\utnrect{ptGray}{0.75}{0.216}{0.196}{0.632}{0.128}
\utncode[style=css,x=0.216,y=0.346,w=0.320,h=0.333,bg=FFFFFF,fs=6.29,lh=8.87,padx=2.0,pady=0.5]{css-div.css}
\utnrect{ptGray}{0.75}{0.216}{0.346}{0.320}{0.333}
\utnpic[height=0.158\paperheight,keepaspectratio=false]{div-resultado}{0.105}{0.767}{0.815}
\utnrect{ptGray}{0.75}{0.105}{0.767}{0.815}{0.158}
\end{frame}

% 56
\begin{frame}{Modelo de cajas - Fondo}\relax
\begingroup
\setbeamerfont{itemize/enumerate body}{size=\fontsize{11.9}{14.7}\selectfont}
\setbeamerfont{itemize/enumerate subbody}{size=\fontsize{11.0}{13.6}\selectfont}
\setbeamerfont{itemize/enumerate subsubbody}{size=\fontsize{10.1}{12.5}\selectfont}
\begin{itemize}
  \item\relax El fondo puede ser un color simple o una imagen. Solamente se visualiza en el área ocupada por el contenido y su relleno.
  \item\relax CSS tiene 5 propiedades para definir el fondo de cada elemento:
  \begin{itemize}
    \item\relax background-color
    \item\relax background-image
    \item\relax background-repeat\flechag{} \textit{{\fontsize{9.1}{11.9}\selectfont si la imagen que queremos mostrar es más pequeña que la caja, entonces la misma por defecto se repetirá horizontal y verticalmente. Con este atributo, podemos cambiar el comportamiento por defecto, e indicar que solo lo haga en un sentido o que no se repita.}}
    \item\relax background-attachment
    \item\relax background-position \flechag{} \textit{{\fontsize{9.1}{11.9}\selectfont permite controlar la posición de la imagen dentro del fondo del elemento mediante la propiedad}}
  \end{itemize}
\end{itemize}
\endgroup
\end{frame}

% 57
\begin{frame}{Posicionamiento}\relax
\begingroup
\setbeamerfont{itemize/enumerate body}{size=\fontsize{13.2}{18.0}\selectfont}
\setbeamerfont{itemize/enumerate subbody}{size=\fontsize{12.2}{16.6}\selectfont}
\setbeamerfont{itemize/enumerate subsubbody}{size=\fontsize{11.2}{15.2}\selectfont}
\begin{itemize}
  \item\relax En HTML los elementos se pueden agrupar en dos grandes categorías
  \begin{itemize}
    \item\relax \textbf{block}: siempre empiezan en una nueva línea y ocupan todo el espacio disponible hasta el final de la línea
    \item\relax \textbf{inline}: no empiezan \textbf{necesariamente} en nueva línea y sólo ocupan el espacio necesario para mostrar sus contenidos.
  \end{itemize}
  \item\relax Esto tiene incidencia en el tamaño de las cajas
\end{itemize}
\endgroup
\utnpic[height=0.280\paperheight,keepaspectratio=false]{block-inline}{0.170}{0.671}{0.653}
\utnrect{black}{0.75}{0.170}{0.671}{0.653}{0.280}
\end{frame}

% 58
\begin{frame}{Posicionamiento}\relax
\begingroup
\setbeamerfont{itemize/enumerate body}{size=\fontsize{12.2}{15.1}\selectfont}
\setbeamerfont{itemize/enumerate subbody}{size=\fontsize{10.2}{12.6}\selectfont}
\setbeamerfont{itemize/enumerate subsubbody}{size=\fontsize{9.3}{11.6}\selectfont}
\begin{itemize}
  \item\relax CSS define cinco modelos diferentes para posicionar una caja:
  \begin{itemize}
    \item\relax \textbf{estático}: se usa por defecto en navegadores.
    \item\relax \textbf{relativo}: posiciona una caja según el posicionamiento estático y después desplazarla respecto de su posición original.
    \item\relax \textbf{absoluto}: la posición de una caja se establece de forma absoluta respecto de su elemento contenedor.
    \item\relax \textbf{fijo}: convierte una caja en un elemento inamovible. Su posición en la pantalla siempre es la misma independientemente del resto de elementos o de si el usuario sube o baja la página en el navegador.
    \item\relax \textbf{flotante}: desplaza las cajas todo lo posible hacia la izquierda o hacia la derecha de la línea en la que se encuentran.
  \end{itemize}
  \item\relax El atributo “position” define el tipo de posicionamiento y los atributos top, right, bottom y ~left indican el desplazamiento
\end{itemize}
\endgroup
\end{frame}

% 59
\begin{frame}{Posicionamiento}\relax
\begingroup
\setbeamerfont{itemize/enumerate body}{size=\fontsize{12.7}{17.3}\selectfont}
\setbeamerfont{itemize/enumerate subbody}{size=\fontsize{11.7}{15.9}\selectfont}
\setbeamerfont{itemize/enumerate subsubbody}{size=\fontsize{10.7}{14.6}\selectfont}
\begin{itemize}
  \item\relax El posicionamiento estático sólo se tiene en cuenta si el elemento es de bloque o en línea, sus propiedades width y height y su contenido.
  \item\relax El posicionamiento relativo desplaza una caja respecto de su posición el posicionamiento estático. El desplazamiento de la caja se controla con las propiedades top, right, bottom y left.
\end{itemize}
\endgroup
\utnpic[height=0.298\paperheight,keepaspectratio=false]{posicionamiento}{0.184}{0.662}{0.632}
\end{frame}

% 60
\begin{frame}{Visualización}\relax
\begingroup
\setbeamerfont{itemize/enumerate body}{size=\fontsize{12.2}{15.1}\selectfont}
\setbeamerfont{itemize/enumerate subbody}{size=\fontsize{11.3}{14.0}\selectfont}
\setbeamerfont{itemize/enumerate subsubbody}{size=\fontsize{10.4}{12.9}\selectfont}
\begin{itemize}
  \item\relax Con las propiedades display y visibility ~se controla la visualización de los elementos. 
  \begin{itemize}
    \item\relax permiten ocultar cualquier elemento de la página. 
    \item\relax se utilizan junto con JavaScript para crear efectos dinámicos.
  \end{itemize}
  \item\relax “\textbf{display}” permite ocultar completamente un elemento haciendo que desaparezca de la página. 
  \begin{itemize}
    \item\relax Como el elemento oculto no se muestra, el resto de elementos de la página se mueven para ocupar su lugar.
  \end{itemize}
  \item\relax “\textbf{visibility}” hace invisible un elemento, lo que significa que el navegador crea la caja del elemento pero no la muestra. 
  \begin{itemize}
    \item\relax El resto de elementos de la página no modifican su posición, ya que aunque la caja no se ve, sigue ocupando sitio.
  \end{itemize}
\end{itemize}
\endgroup
\end{frame}

% 61
\begin{frame}{Visualización}\relax
\begingroup
\setbeamerfont{itemize/enumerate body}{size=\fontsize{12.2}{16.6}\selectfont}
\setbeamerfont{itemize/enumerate subbody}{size=\fontsize{11.2}{15.2}\selectfont}
\setbeamerfont{itemize/enumerate subsubbody}{size=\fontsize{10.3}{14.0}\selectfont}
\begin{itemize}
  \item\relax Ejemplo donde se nota la diferencia entre display (que deja libre el lugar cuando no está seteado) y visibility que simplemente no dibuja el elemento pero deja el espacio reservado.
\end{itemize}
\endgroup
\utnpic[height=0.505\paperheight,keepaspectratio=false]{display-visibility}{0.227}{0.406}{0.724}
\end{frame}

% 62
\begin{frame}{Texto}\relax
\begingroup
\setbeamerfont{itemize/enumerate body}{size=\fontsize{12.2}{14.6}\selectfont}
\setbeamerfont{itemize/enumerate subbody}{size=\fontsize{10.2}{12.2}\selectfont}
\setbeamerfont{itemize/enumerate subsubbody}{size=\fontsize{9.3}{11.2}\selectfont}
\begin{itemize}
  \item\relax Para modificar el aspecto de un texto tenemos varios atributos:
  \begin{itemize}
    \item\relax \textbf{color}: define el color del texto
    \item\relax \textbf{font-family}: se puede indicar de dos maneras, con el nombre de una tipográfía ("Arial", "Verdana", "Garamond“) o con el nombre genérico de una familia tipográfía que hace referencia al estilo del tipo de letra ( serif - sans-serif - cursive, fantasy y monospace).
    \item\relax \textbf{font-size}: establece el tamaño de letra utilizado para el texto. Puede indicarse como una medida de tamaño: ~absoluto (xx-small, x-small, small, medium, large, x-large, xx-large) ~- relativo ( solo se indica large o smaller y será en relación al elemento padreo)- unidad de medida – porcentaje.
    \item\relax \textbf{font-weight}: establece el ancho de la fuente. Puede tomar valores normal | bold | bolder | lighter
    \item\relax \textbf{font-style}: stilo de la letra que puede ser normal | italic | oblique
  \end{itemize}
\end{itemize}
\endgroup
\end{frame}

% 63
\begin{frame}{Texto}\relax
\begingroup
\setbeamerfont{itemize/enumerate body}{size=\fontsize{12.2}{14.6}\selectfont}
\setbeamerfont{itemize/enumerate subbody}{size=\fontsize{10.2}{12.2}\selectfont}
\setbeamerfont{itemize/enumerate subsubbody}{size=\fontsize{9.3}{11.2}\selectfont}
\begin{itemize}
  \item\relax También se pueden modificar aspectos mas elaborados
  \begin{itemize}
    \item\relax \textbf{text-align}: indica la alineación del texto si será left | right | center | justify
    \item\relax \textbf{line-height}: indica el espacio entre las líneas de un texto
    \item\relax \textbf{text-decoration}: indica la decoración que tendrá, los valores pueden ser: none - underline -overline - line-through - blink (parpadeante)
    \item\relax \textbf{text-transform}: transforma las letras del texto original, en mayúsculas o minúsculas (capitalize - uppercase - lowercase- none)
    \item\relax \textbf{white-space}: indica como tratar a los espacios en blanco. (normal | pre | nowrap | pre-wrap | pre-line )
    \item\relax \textbf{word-spacing} y \textbf{letter-spacing}: permiten establecer el espacio entre las letras o las palabras.
    \item\relax \textbf{text-indent}: Tabula desde la izquierda la primera línea del texto original la unidad de medida indicada.
  \end{itemize}
\end{itemize}
\endgroup
\utnpic[height=0.263\paperheight,keepaspectratio=false]{meme-bigote}{0.590}{0.672}{0.254}
\end{frame}

% 64
\begin{frame}{Enlaces}\relax
\begingroup
\setbeamerfont{itemize/enumerate body}{size=\fontsize{12.2}{15.1}\selectfont}
\setbeamerfont{itemize/enumerate subbody}{size=\fontsize{10.2}{12.6}\selectfont}
\setbeamerfont{itemize/enumerate subsubbody}{size=\fontsize{9.3}{11.6}\selectfont}
\begin{itemize}
  \item\relax Al igual que con los textos, con los enlaces podemos modificar el tamaño de letra, su color y la decoración del texto del enlace.
  \item\relax Pero además se pueden aplicar diferentes estilos a un mismo enlace en función de su estado. (visitado, o cuando el mouse pasa por encima por ejemplo).
  \item\relax Para esto se incluyen en CSS “pseudo-clases” que permiten aplicar más de un estilo a un selector según su estado
  \begin{itemize}
    \item\relax :link, aplica estilos a los enlaces que apuntan a páginas o recursos que aún no han sido visitados por el usuario.
    \item\relax :visited, aplica estilos a los enlaces que apuntan a recursos que han sido visitados anteriormente por el usuario
    \item\relax :hover, aplica estilos al enlace sobre el que está el mouse.
    \item\relax :active, aplica estilos al enlace que está haciendo click el usuario.
  \end{itemize}
\end{itemize}
\endgroup
\end{frame}

% 65
{\setbeamertemplate{enumerate item}{\insertenumlabel)}\begin{frame}{Ejercicio 2 (CSS)}\relax
\begingroup
\setbeamerfont{itemize/enumerate body}{size=\fontsize{11.5}{13.7}\selectfont}
\setbeamerfont{itemize/enumerate subbody}{size=\fontsize{10.5}{12.6}\selectfont}
\setbeamerfont{itemize/enumerate subsubbody}{size=\fontsize{9.7}{11.6}\selectfont}
{\fontsize{11.5}{13.7}\selectfont Añadir \textbf{estilos} a la “Página de DesoTravel”, conectando un archivo \textbf{.css} con el \textbf{.html} anterior.\par}
\begin{itemize}
  \item\relax Los estilos son libres, pero deben cumplir las siguientes condiciones:
\end{itemize}
\begin{enumerate}
  \item\relax Utilizar al menos un <section> o un <article>
  \item\relax Utilizar selectores de elemento, ID y de clase
  \item\relax Centrar el título principal y la foto
  \item\relax Colorear el fondo, ya sea con una imagen o con un color sólido (también se pueden hacer gradientes): \href{https://cssgradient.io/}{\uline{https:/\allowbreak{}/\allowbreak{}cssgradient.io/\allowbreak{}}} 
  \item\relax Separar las partes de “Servicios disponibles” (color verde / más amigable) de las “Condiciones de cancelación” (color rojo / menos amigable).
\end{enumerate}
\endgroup
\end{frame}
}

% 66
\begin{frame}{Flexbox Layout}\relax
\begingroup
\setbeamerfont{itemize/enumerate body}{size=\fontsize{12.2}{16.6}\selectfont}
\setbeamerfont{itemize/enumerate subbody}{size=\fontsize{11.2}{15.2}\selectfont}
\setbeamerfont{itemize/enumerate subsubbody}{size=\fontsize{10.3}{14.0}\selectfont}
\begin{itemize}
  \item\relax Existen múltiples layouts como son \textbf{flexbox (unidimensional)} y \textbf{grid (bidimensional)}
  \item\relax Básicamente nos ayudan a crear layouts dinámicos y flexibles de una manera sencilla
  \item\relax Se debe definir el contenedor padre con \textbf{display: flex;} y desde ahí se puede direccionar con el atributo \textbf{flex-direction}
  \item\relax Más información interesante: \href{https://css-tricks.com/snippets/css/a-guide-to-flexbox/}{\uline{https:/\allowbreak{}/\allowbreak{}css-tricks.com/\allowbreak{}snippets/\allowbreak{}css/\allowbreak{}a-guide-to-flexbox/\allowbreak{}}} 
\end{itemize}
\endgroup
\utnpic[height=0.299\paperheight,keepaspectratio=false]{meme-backend}{0.611}{0.701}{0.166}
\end{frame}

\end{document}
